%% ============================================================
%% PAPER 1: TOPOLOGY ABSTRACTION & PHYSICS FIDELITY EFFECTS
%% Target Journal: Applied Energy
%% Version: 14.1 — Terminology patch: "synthetic sweep" removed
%% ============================================================

\documentclass[a4paper,fleqn]{cas-dc}

%% --- Bibliographie ---
\usepackage[numbers,sort&compress]{natbib}

%% --- Mathematik ---
\usepackage{amsmath,amssymb,amsfonts}

%% --- Tabellen ---
\usepackage{booktabs}
\usepackage{multirow}
\usepackage{array}
\usepackage{tabularx}
\usepackage{threeparttable}

%% --- Grafik & Floats ---
\usepackage{graphicx}
\usepackage{float}
\usepackage{subfig}
\usepackage[section]{placeins}

%% --- Einheiten & Chemie ---
\usepackage{siunitx}
\usepackage[version=4]{mhchem}

%% --- Sonstiges ---
\usepackage{xcolor}
\usepackage{enumitem}
\usepackage{pifont}
\usepackage{adjustbox}

%% --- Hyperlinks (immer spät laden) ---
\usepackage{hyperref}
\usepackage{cleveref}
\usepackage[numbers]{natbib}
%% --- SI-Setup ---
\sisetup{
  per-mode=symbol,
  range-phrase=--,
  range-units=single
}

%% --- Hyperref-Setup ---
\hypersetup{
  colorlinks=true,
  linkcolor=blue,
  citecolor=blue,
  urlcolor=blue
}

%% --- Float-Platzierung maximal lockern ---
\renewcommand{\topfraction}{0.95}
\renewcommand{\bottomfraction}{0.95}
\renewcommand{\textfraction}{0.05}
\renewcommand{\floatpagefraction}{0.85}
\setcounter{topnumber}{4}
\setcounter{bottomnumber}{4}
\setcounter{totalnumber}{8}

%% --- Speziell für double-column (table*/figure*) ---
\renewcommand{\dbltopfraction}{0.95}
\renewcommand{\dblfloatpagefraction}{0.85}

%% --- Seiten dürfen unten unterschiedlich lang sein ---
\raggedbottom

%% --- Custom commands ---
\newcommand{\placeholder}[1]{\textcolor{red}{[#1]}}
\newcommand{\Ltwo}{\mathrm{L2}}
\newcommand{\Lthree}{\mathrm{L3}}
\newcommand{\Lplus}{\mathrm{L3}^{+}}
\newcommand{\Lnl}{\mathrm{L3}^{\mathrm{NL}}}

%% --- Einheitliche Formel-Formatierung ---
\setlength{\mathindent}{1em}
\setlength{\abovedisplayskip}{8pt plus 2pt minus 4pt}
\setlength{\belowdisplayskip}{8pt plus 2pt minus 4pt}
\setlength{\abovedisplayshortskip}{4pt plus 2pt}
\setlength{\belowdisplayshortskip}{6pt plus 2pt minus 2pt}

%% --- Custom column types ---
\newcolumntype{P}[1]{>{\centering\arraybackslash}m{#1}}
\newcolumntype{Q}[1]{>{\raggedright\arraybackslash}m{#1}}

%% --- Fix caption alignment for cas-dc ---
\makeatletter
\long\def\@makecaption#1#2{%
  \vskip\abovecaptionskip
  \raggedright
  {\bfseries #1.}\enspace #2\par
  \vskip\belowcaptionskip}
\makeatother

\begin{document}
\let\WriteBookmarks\relax
\def\floatpagepagefraction{1}
\def\textpagefraction{.001}

\shorttitle{Topology and Physics Fidelity in District Heating Dispatch}
\shortauthors{L.~Ruess and A.~Sauer}

\title[mode = title]{Topology Resolution Dominates Dispatch Accuracy in District Heating Networks with Industrial Demands: A Five-Level MILP Comparison}

\author[1,2]{Lukas Ruess}[orcid=0009-0005-5010-3520]
\cormark[1]
\ead{lukas.ruess@ipa.fraunhofer.de}
\credit{Conceptualization, Methodology, Software, Validation,
        Formal analysis, Investigation, Data curation,
        Writing -- original draft, Visualization, Funding acquisition}



\author[1,2]{Alexander Sauer}
\credit{Supervision, Writing -- review \& editing}


\affiliation[1]{organization={Fraunhofer Institute for Manufacturing
                Engineering and Automation IPA},
                addressline={Nobelstr.~12},
                city={Stuttgart},
                postcode={70569},
                country={Germany}}
\affiliation[2]{organization={Institute for Energy Efficiency in
                Production EEP, University of Stuttgart},
                addressline={Allmandring 35},
                city={Stuttgart},
                postcode={70569},
                country={Germany}}



\cortext[1]{Corresponding author.}

\begin{abstract}
The widespread assumption that higher physics fidelity yields
proportional accuracy gains in district heating dispatch is
untested. This study provides a comparison across
184 optimisation  instances on a validated industrial network and
36~synthetic topologies, pressure-drop hydraulics, temperature
propagation, and nonlinear formulations collectively contribute
less than 1\,\% to dispatch cost while the computationally
trivial step of resolving network topology captures
13--39\,\% of hidden cost bias.

A five-level taxonomy (L1 copperplate through $\Lnl$ MIQCP)
isolates three effects through pairwise comparison. Results show
that topology dominates: the copperplate to detailed gap reaches
13\,\% cost underestimation and 74\,t\,\ce{CO2}/year on the
primary case, driven almost entirely by thermal losses made
visible through spatial resolution. Extended hydraulic physics
add ${\leq}\,0.3\,\%$; the combined linearisation and transport-delay effect remains at or below solver tolerance (${\leq}\,0.5\,\%$) across two independent winter windows,
confirming that MILP linearisation is adequate for planning-accuracy dispatch. These findings yield quantitative
selection thresholds linking pipe length and demand heterogeneity
to the minimum warranted model fidelity.

Within the tested scope (radial trees, fixed heating curve,
dispatch with given capacities), practical implications could be determined: resolve the network graph before refining the physics.
Ring topologies, variable supply temperature, and joint
dispatch--sizing problems may shift these thresholds and are
identified as priority follow-up studies.
\end{abstract}

\begin{keywords}
District heating \sep
Mixed-integer linear programming \sep
Mixed-integer quadratically constrained programming \sep
Topology abstraction \sep
Thermal storage dispatch \sep
Heat pump optimisation  \sep
Network modelling fidelity \sep
Linearisation error
\end{keywords}

\maketitle



%% ============================================================
%% 1. INTRODUCTION
%% ============================================================
\section{Introduction}
\label{sec:introduction}

%\subsection{Motivation and problem statement}
%\label{sec:motivation}

District heating networks (DHN) for industrial usage are central to European decarbonisation
strategies, integrating heat pumps (HP), combined heat and power (CHP), and thermal storage (TES) within
multi-source systems~\cite{Lund2014_4GDH,Werner2017,Lund2025_clean}. A recent bibliometric
analysis confirms that DH research output has grown steadily over
the past decade, with optimisation  and modelling now constituting
one of the field's dominant thematic clusters~\cite{PereaMoreno2026}.
Operational optimisation  of such networks increasingly relies on
mathematical programming. The level of detail in the modelling of heating networks depends heavily on whether the topology and thermo-hydraulics are simplified:

%
\begin{enumerate}[nosep]
  \item \textbf{Topology abstraction:} the level of spatial detail
        at which the network is represented from a single-bus
        copperplate to a pipe-by-pipe graph.
  \item \textbf{Thermo-hydraulic accuracy:} the set of physical phenomena
        modelled from zero loss energy balances to
        pressure drop coupled, temperature propagating, delay aware
        thermo-hydraulic models.
\end{enumerate}

At one extreme, \emph{copperplate} models aggregate all components
and demands to a single bus, ignoring spatial structure and thermal
losses entirely. \citet{Wirtz2020} explicitly acknowledge
this limitation: ``the validity of the global energy balance
assumption depends on the network topology, size and the operation
strategy,'' yet no quantification of the resulting error has been
published. At the other extreme, \emph{full thermo-hydraulic} models
resolve the physical pipe-by-pipe network, capturing thermal losses
 5--15\,\% of annual heat
supply~\cite{Frederiksen2013,Talebi2016}, pressure-dependent
pumping costs, temperature propagation, and transport delays arising
from finite fluid velocity. A further complication arises from the
\emph{mathematical formulation}: extended thermo-hydraulic physics
introduce nonlinear and bilinear terms incompatible with standard
Mixed integer linear programming (MILP) solvers, forcing a choice between linearisation (preserving
MILP tractability but introducing approximation error) and
nonlinear/quadratic formulations (preserving physical accuracy but
potentially sacrificing solve time).

A common assumption is that higher thermo-hydraulic accuracy and finer spatial resolution both yield proportional accuracy
gains, which can justify the associated computational cost. Yet no quantitative evidence exists to confirm or refute this
proportionality, nor to identify which modelling dimension delivers
the dominant return on complexity. A closely related effect has
been quantified for electricity transmission systems: low
spatial network resolution causes capacity-expansion models to
ignore intra-regional congestion and underestimate system costs,
whereas spatial aggregation choices dominate other modelling
simplifications in scenario comparisons~\cite{Frysztacki2021,
FrewJacobson2016}. Whether an analogous dominance of topology over
thermo-hydraulic detail holds for DH dispatch optimisation  has not been established.

\subsection{Research gap}
\label{sec:gap}

The existing literature focuses on studies involving multidimensional adjustments: published comparisons change
multiple modelling dimensions simultaneously. For instance,
comparing the full nonlinear pipe model with temperature-dependent
Coefficient of performance (COP) used by \citet{Vesterlund2017} against the
copperplate formulation of \citet{Gabrielli2018}
conflates topology effects with (i)~None Linear Programming (NLP) vs.\ constant-loss
formulations, (ii)~temperature-dependent vs.\ constant COP,
(iii)~the presence or absence of hydraulic coupling, and
(iv)~differing temporal resolutions. The same pattern appears
systematically across the literature
(\Cref{sec:related_work}). Four specific gaps persist:
%
\begin{enumerate}[nosep]
  \item The isolated topology abstraction as the sole
        experimental variable while holding physics constant.
  \item The isolated thermo-hydraulic effect (adding
        pressure drop, temperature propagation, transport delay)
        with a similar topology.
  \item The linearisation error introduced by MILP-compatible
        approximations of thermo-hydraulic physics has not been
        bounded against a quadratic reference in a DH dispatch
        context.
  \item A quantitative guide to selecting piping topology and
        thermo-hydraulic design based on observable network characteristics.
\end{enumerate}

\subsection{Contributions and research questions}
\label{sec:contributions}

This paper addresses these four gaps through a controlled
experimental design with three orthogonal comparison dimensions.
The controlled comparison simultaneously holds topology
constant while varying physics, and holds thermo-hydraulics constant while
varying topology. This paper makes four contributions, each answering one
research question:
%
\begin{description}[style=nextline, nosep, leftmargin=1.5em]
  \item[Contribution 1 / RQ1 Topology effect:]
    A five-level taxonomy and topology comparison within a unified MILP
    framework where scenarios share thermo-hydraulics. This
    isolates topology as the sole experimental variable and is
    validated against measured industrial operational data.
    \emph{How much does topology abstraction in the dispatch optimisation  change annual
    operational cost and \ce{CO2} emissions?}

  \item[Contribution 2 / RQ2 thermo-hydraulic effect:]
    A controlled thermo-hydraulic accuracy comparison
    that isolates the effect of adding pressure drop and
    temperature-dependent loss propagation while holding topology
    and formulation class constant.
    \emph{Do advanced hydraulic and thermal calculations
    have a significant impact on load distribution in the dispatch optimisation ?}

  \item[Contribution 3 / RQ3 Linearisation error bounds:]
    An upper-bound characterisation of the combined linearisation
    and transport-delay effect against an mixed integer quadratic constraint programming (MIQCP) reference
    ($\Lplus$ vs.\ $\Lnl$).
    \emph{What is the maximum cost deviation attributable to MILP
    linearisation and transport-delay omission?}

  \item[Contribution 4 / RQ4 Generalisability:]
    Parametric analysis across 36~synthetic configurations
    yielding joint selection rules based on observable network
    properties.
    \emph{Which network properties determine whether detailed
    topology and/or extended physics representation is warranted in dispatch optimisation ?}
\end{description}

%% ============================================================
%% 2. RELATED WORK
%% ============================================================
\section{Related work}
\label{sec:related_work}

In the field of dispatch strategies for DHN numerous studies have been published. Given the wide range of applications, MILP modelling strategies for the topologies of DHN are a frequent topic of discussion (\Cref{sec:rw_milp_topo}). The linearisation of thermo-hydraulic networks varies as well across these studies (\Cref{sec:rw_thermo-hydraulic}).

\subsection{MILP optimisation  and spatial scale in DH systems}
\label{sec:rw_milp_topo}

MILP has become the dominant methodology for operational
optimisation  of DH systems~\cite{Fang2016,Guelpa2019review}.
Open-source frameworks like oemof~\cite{Hilpert2018_oemof},
PyPSA~\cite{Brown2018_pypsa}, FINE~\cite{Welder2018_FINE},
Calliope~\cite{Pfenninger2018_calliope},
urbs~\cite{Dorfner2016_urbs} support multi-carrier energy
system optimisation  but share common limitations for detailed DH
analysis: thermal losses are typically time-invariant, heat pump
COP is constant or seasonal, hydraulic coupling is absent, and
network topology is either fully aggregated or fixed at a single
detail level. All five frameworks operate in the default setting at L1 in the taxonomy proposed here; the cost of this simplification has not been
quantified.

The DH modelling literature exhibits a continuous spectrum of
spatial aggregation. The reduction of detailed pipe networks to a
small number of aggregated branches has a long tradition in DH
simulation: \citet{Larsen2004} compare aggregation methods on a
23-substation network and show that pipe count can be reduced by
roughly an order of magnitude before production and return-temperature
errors become significant, while \citet{Falay2020} extend this line
of work to large-scale dynamic simulation of several thousand
substations. \citet{Schweiger2018_paradigms} survey the resulting
landscape of DH modelling paradigms and general-purpose tools. These
studies establish \emph{which} aggregation preserves simulation
fidelity, but none quantifies the resulting bias in
\emph{dispatch-optimisation  cost}.
Single-node models are standard in
national-scale analyses~\cite{Connolly2014_EnergyPLAN,
Brown2018_pypsa,Lund2017_SmartEnergy}.
\citet{Gabrielli2018} formulated a MILP for distributed
multi-energy systems including thermal and electrical storage but
treat the thermal network as a single aggregated efficiency
factor. Zone-based models (3--15~nodes) aggregate demand into
geographic clusters~\cite{Wirtz2020,Mertz2016,Buffa2019};
\citet{Buenning2020} demonstrate that zone-level aggregation
captures 85--95\,\% of dynamic temperature variation. Full-graph
models resolve node-by-node from GIS
data~\cite{Vesterlund2017,Schweiger2017,vanderHeijde2017,
Blommaert2020}; \citet{Giraud2015} show tractability up to
${\approx}\,100$~nodes for hourly annual horizons.

Across this spectrum, the level of spatial aggregation has been
shown to materially affect dispatch
results~\cite{Wirtz2020,Buffa2019}, but no scalar metric links
observable network properties to the required resolution level.
\citet{Leitner2019} compared copperplate and full-graph
representations for an Austrian DH network and found cost
differences of 3--8\,\%, but simultaneously changed the loss (more instances in \Cref{sec:rw_summary}).
Three very recent studies extend this line without isolating
topology as an independent factor. \citet{Trentmann2026} integrate
full-graph DHN expansion into an energy-hub MILP with
temporally-resolved heat-pump COPs, showing that network supply
temperature strongly shifts the optimal centralized-vs-building-level
supply split; their model always resolves the network at
full-graph detail, so the copperplate/aggregated comparison studied
here is outside their scope. \citet{Ceruti2025} likewise formulate a
full-graph MILP for DHN design with TES and find that temporal
(not spatial) resolution dominates design quality. \citet{Vieth2026} address expansion
planning of existing full-graph networks under resilience
constraints, showing that considering generation-unit outages
lowers the connectable peak demand by up to 97\,\%. 

A parallel and largely separate literature quantifies the dispatch
value of the generation technologies studied in the present case
(HP, electrode boiler, CHP, TES) without addressing spatial
resolution. \citet{Nielsen2016} and \citet{Boettger2015} show, for
the Danish and German electricity markets respectively, that
deterministic dispatch models systematically overestimate the
economic value of HP and electric-boiler flexibility relative to
stochastic or control-power-aware formulations. \citet{Galteland2026} confirm this technology-value
question is still actively studied: for a hybrid industrial
electrode-boiler/steam-accumulator/battery system, they show optimal
sizing is highly sensitive to reserve-market participation and
storage cost assumptions. \citet{Mollenhauer2018} formulate an MILP
unit-commitment and dispatch problem for German CHP plants
retrofitted with HP and TES at hourly resolution over a full year,
methodologically close to the present dispatch model but, like
\citet{Schweiger2017_p2h,Averfalk2017,David2017},
representing the DH network itself as a single aggregated demand
node. \citet{Walden2023} extend this line to a nonlinear industrial
power-to-heat system coupling a high-temperature heat pump, sensible
TES, and wind generation, optimizing operation with an accurate
nonlinear component model rather than a network representation;
their setting shares this paper's technology mix (HP-based
electrification, TES buffering of variable renewable supply) but,
like \citet{Mollenhauer2018}, treats the distribution network as a
single lumped point rather than a resolved topology.
\citet{Riedl2026} propose, most recently, a general flexibility
assessment framework for DHC systems that explicitly incorporates
network-side constraints alongside generation-side flexibility , though their published
case studies do not yet vary spatial resolution as an experimental
factor. \citet{Buhler2017} instead quantify the industrial excess-heat
resource available to DH networks spatially, but not within a
dispatch-optimisation  framework. None of these studies varies
network topology as an experimental factor, reinforcing that the
topology/technology interaction remains unexplored.

\subsection{Thermo-hydraulic network models and linearisation}
\label{sec:rw_thermo-hydraulic}

Beyond topology, the thermo-hydraulic fidelity of pipe models varies
substantially. Steady-state thermal losses follow from the
fluid--ground temperature difference and pipe insulation
properties and are MILP-compatible when
temperatures are fixed parameters ~\cite{Frederiksen2013}. The Darcy--Weisbach equation,
in contrast, relates pressure drop to the square of flow velocity,
introducing a quadratic nonlinearity; pumping power is proportional
to $\dot{m} \cdot \Delta p$, creating a cubic term. Authors have
addressed this through piecewise
linearisation~\cite{Bordin2016,Mertz2016} or second-order cone
relaxations~\cite{Blommaert2020}. \citet{Hari2024} provide a
comprehensive thermal network flow optimisation  framework
integrating Darcy--Weisbach pressure drop with explicit pump
placement. Reported pumping costs range from 1--5\,\% of thermal
energy costs~\cite{Frederiksen2013}.

Temperature propagation along pipes introduces exponential and
bilinear terms~\cite{vanderHeijde2017}; \citet{Maldonado2024}
validate this model against industrial data, achieving supply
temperature Mean absolute error (MAE) below \SI{0.5}{\kelvin}. Transport delays of 1--3\,h arise for
long pipes ~\cite{vanderHeijde2017,Giraud2015}; the
fixed-delay approximation is most accurate when flow variations
are moderate ($\pm 30\,\%$ of nominal).

MIQCP formulations offer a middle ground between MILP and full
NLP~\cite{Ommen2016}: quadratic constraints are handled natively
by modern solvers while integer variables retain unit-commitment
logic. \citet{Hering2021} and \citet{Hering2022} demonstrate this
approach for DH design and dispatch with multiple heat pumps,
formulating temperature-dependent constraints as MIQCP for a 4GDH
network; their case studies show the MIQCP remains tractable at the
scale of tens of consumers, but do not report a reference-model
comparison isolating the cost of PWL-MILP approximation relative to
the native quadratic formulation. However, MIQCP problems lack
LP-relaxation tightness, potentially increasing solve times. The
magnitude of the linearisation error in a DH dispatch context i.e.,
the cost gap between Piece wise linear (PWL)-MILP and native MIQCP
formulations has not been published.

\subsection{Positioning and research gap}
\label{sec:rw_summary}
\Cref{tab:prior_work} reveals the systematic confounding of
eleven representative studies spanning copperplate to full-graph
resolution, \emph{none} isolates topology from physics as an
independent experimental variable.
\Cref{tab:prior_work} positions representative studies. No prior work isolates topology from
physics, nor bounds linearisation error against a quadratic
reference in a controlled DH dispatch setting.

Two studies come closest to the present design. \citet{Quaggiotto2021}
report, for the Verona DH network, that simplifying the network's
thermal behaviour in a dispatch optimisation  underestimates
operational cost relative to a model-predictive controller using a
detailed network model, but without a taxonomy that separates topology from
thermo-hydraulic fidelity or quantifies the gap as a function of
network properties. \citet{Jansen2024} independently confirm this
direction with a mixed-integer non-linear MPC for DH networks,
reporting a similar fidelity/performance trade-off for a real
Belgian network, but again without isolating topology as a factor.
\citet{Wirtz2021_complexity} systematically vary
the level of detail across 24 MILP formulations of a multi-energy
system and identify which simplifications are cost-relevant; their
experimental design closely parallels the controlled-comparison logic
adopted here, but their model complexity axis spans technology and
part-load representation rather than network topology and
thermo-hydraulics, and the building-scale multi-energy system they
study does not resolve pipe-by-pipe network structure. \citet{Kotzur2021}
review complexity-management methods for energy system optimisation 
more broadly and note that, while temporal aggregation is
well-studied, spatial aggregation "has only a small community" and
lacks systematic guidance. The present
paper closes the gap by comparing
systematic, axis-isolated logic specifically to the topology/physics
trade-off in DH dispatch, and within a single
framework spanning all three dimensions (topology, thermo-hydraulics,
linearisation).


%% ============================================================
%% 3. METHODOLOGY
%% ============================================================
%% ============================================================
%% 3. METHODOLOGY
%% ============================================================
\section{Methodology}
\label{sec:methodology}

This section defines the controlled experimental design
(\Cref{sec:exp_design}), presents the mathematical formulations in
order of increasing fidelity (\Cref{sec:milp,sec:extended}),
describes the validation protocol (\Cref{sec:val_protocol}), and
closes with implementation details (\Cref{sec:model_impl}).

\begin{table*}[!tbp]
\centering
\caption{DH optimisation  studies with
isolated topology and thermo-hydraulics as experimental
variables.}
\label{tab:prior_work}
\begin{threeparttable}
\footnotesize
\setlength{\tabcolsep}{4pt}
\begin{tabular}{@{}l c c c c c c c@{}}
\toprule
Reference & Topo. & Physics & Form. & COP &
  Scale & Temp. & Comparison \\
\midrule
\citet{Connolly2014_EnergyPLAN}
  & Cop. & Fix.~$\eta$ & LP & Const. & Nat. & Ann. & No \\
\citet{Gabrielli2018}
  & Cop. & Fix.~$\eta$ & MILP & Const. & 10\,MW & Ann. & No \\
\citet{Wirtz2020}
  & Zone(5) & Lin.\ loss & MILP & Seas. & 25\,MW & Ann. & No \\
\citet{Mertz2016}
  & Zone(8) & Lin.+$\Delta p$ & MINLP & Const.
  & 20\,MW & Ann. & No \\
\citet{Hari2024}
  & Full & DW+pumps & NLP & -- & 50\,MW & Steady & No \\
\citet{Leitner2019}
  & L1+L3 & Diff.\ phys. & MILP & Const. & 12\,MW
  & Ann. & No \\
\citet{Bordin2016}
  & Full & PWL~$\Delta p$ & MILP & -- & 40\,MW & Ann. & No \\
\citet{Giraud2015}
  & Full & Steady-st. & NLP & $T$-dep. & 30\,MW & Ann. & No \\
\citet{Vesterlund2017}
  & Full & NLP therm. & NLP & $T$-dep. & 50\,MW & Ann. & No \\
\citet{vanderHeijde2017}
  & Full & Pipe comp. & Var. & -- & Review & --
  & Part.\tnote{a} \\
\citet{Blommaert2020}
  & Full & Adj.\ NLP & NLP & $T$-dep. & 30\,MW & Des. & No \\
\citet{Quaggiotto2021}
  & Full & FD therm. & MPC & -- & Dist. & Days & No \\
\citet{Jansen2024}
  & Full & MINLP & MPC & -- & Dist. & Days & No \\
\citet{Hering2021}
  & Fix. & MIQCP & MIQCP & $T$-dep. & 7 cons. & Des. & No \\
\citet{Wirtz2021_complexity}
  & Cop. & 24 var. & MILP & Var. & Bldg. & Ann. & Part.\tnote{b} \\
\midrule
\textbf{This work}
  & \textbf{$L1$--$\Lnl$}
  & Bas.+ext.
  & MILP+QCP
  & Pre-c.
  & 12\,MW
  & 8760\,h
  & Yes \\
\bottomrule
\end{tabular}
\begin{tablenotes}\footnotesize
  \item[a] Compared pipe-loss formulations but not spatial
    aggregation levels.
  \item[b] Systematically varied model complexity (24 variants)
    but along technology/part-load axes, not network topology.
\end{tablenotes}
\end{threeparttable}
\end{table*}

%% ============================================================
%% 3.1 EXPERIMENTAL DESIGN
%% ============================================================
\subsection{Controlled experimental design}
\label{sec:exp_design}

The experimental design isolates three effects through pairwise
comparisons (\Cref{fig:comparison_design}):
%
\begin{enumerate}[nosep]
  \item \textbf{Topology effect} (L1\,$\to$\,L2\,$\to$\,L3):
        spatial resolution varies; physics held constant.
  \item \textbf{Physics fidelity effect} (L3\,$\to$\,$\Lplus$):
        physics scope varies; topology and formulation class (MILP)
        held constant.
  \item \textbf{Linearisation error} ($\Lplus$\,$\to$\,$\Lnl$):
        mathematical formulation varies (MILP\,$\to$\,MIQCP);
        physics and topology held identical.
\end{enumerate}

\begin{figure}[ht]
  \centering
  \includegraphics[width=\linewidth]{paper1_dh_fidelity/F1_experimental_design.png}
  \caption{Experimental design with three orthogonal comparison
  dimensions. Each arrow isolates exactly one modelling dimension.}
  \label{fig:comparison_design}
\end{figure}

\Cref{tab:model_scope} summarises which physical phenomena enter at
each level. The five variants are built up as follows, inspired by \citet{Blommaert2020}.

\begin{table}[htbp]
\centering
\caption{Model scope across all five variants. ``PWL'': piecewise-linear;
``Quad./Bilin.'': native quadratic/bilinear constraint.}
\label{tab:model_scope}
\begin{threeparttable}
\scriptsize
\setlength{\tabcolsep}{2pt}
\begin{tabular*}{\columnwidth}{@{\extracolsep{\fill}}
                    l c c c c c @{}}
\toprule
Phenomenon & L1 & L2 & L3 & $\Lplus$ & $\Lnl$ \\
\midrule
\multicolumn{6}{@{}l@{}}{\textit{Thermal}} \\[1pt]
\quad Pipe losses ($UL\Delta T$)
  & -- & \checkmark & \checkmark & \checkmark & \checkmark \\
\quad Temp.\ propagation
  & -- & -- & -- & PWL & Bilin. \\
\quad Time-varying COP
  & \checkmark & \checkmark & \checkmark & \checkmark & \checkmark \\
\quad Storage losses \& $\eta$
  & \checkmark & \checkmark & \checkmark & \checkmark & \checkmark \\
\midrule
\multicolumn{6}{@{}l@{}}{\textit{Hydraulic}} \\[1pt]
\quad Pressure drop (DW)
  & -- & -- & -- & PWL & Quad. \\
\quad Pumping power
  & -- & -- & -- & \checkmark & \checkmark \\
\midrule
\multicolumn{6}{@{}l@{}}{\textit{Temporal}} \\[1pt]
\quad Transport delay
  & -- & -- & -- & -- & \checkmark \\
\quad Hourly grid $e$-factors
  & \checkmark & \checkmark & \checkmark & \checkmark & \checkmark \\
\midrule
\multicolumn{6}{@{}l@{}}{\textit{Outside the scope of the research}} \\[1pt]
\quad Sub-hourly transients
  & \multicolumn{5}{c}{excluded} \\
\quad Capacity sizing
  & \multicolumn{5}{c}{excluded} \\
\quad $T_{\mathrm{sup}}$ optim.
  & \multicolumn{5}{c}{excluded} \\
\quad Ring / bidir.\ flow
  & \multicolumn{5}{c}{excluded} \\
\bottomrule
\end{tabular*}
\begin{tablenotes}\scriptsize
  \item[] Physics scope as shown applies to the
    \emph{primary case}. In the synthetic parametric study,
    level labels are reassigned to enable the additive
    decomposition; see \Cref{tab:synth_level_mapping}.
\end{tablenotes}
\end{threeparttable}
\end{table}

\paragraph{Level~1 (Copperplate).}
Single bus, no spatial structure, no pipe losses. All assets are
co-located at a single node; the network is represented by a
single scalar heat balance.

\paragraph{Level~2 (Simplified multi-node).}
Aggregated demand zones (7 in the primary case) connected by
simplified pipe segments that preserve the tree topology. The total
$U \cdot L$ product matches L3, ensuring identical aggregate annual
loss energy. Asset locations correspond to their physical positions
(\Cref{sec:case_studies}).

\paragraph{Level~3 (Detailed multi-node, basic MILP).}
Full spatial resolution (15 junction nodes). Individual pipe
parameters with steady-state losses, fixed supply temperature
(heating curve), and pre-computed COP (\Cref{app:cop}). No pressure
drop, no temperature propagation, no transport delay.

\paragraph{$\Lplus$ (Detailed multi-node, extended MILP).}
Same topology as L3. Adds piecewise-linear pressure drop
(Darcy--Weisbach), linearised temperature propagation, and the
resulting in the pumping power. Transport delay is configured but bypassed:
the linearised pipe model cannot represent the source-to-consumer
lag, so this effect is characterised exclusively through the
$\Lplus\!\to\!\Lnl$ comparison.

\paragraph{$\Lnl$ (Detailed multi-node, extended MIQCP).}
Same topology and physics scope as $\Lplus$, but pressure drop
enters as a native quadratic constraint, temperature propagation as
bilinear products, and transport delay is dispatched natively.
$\Lnl$ is the only variant in which transport delay is active.

\paragraph{Auxiliary variant: L1$_{\mathrm{topo}}$.}
For the additive gap decomposition (\Cref{sec:case_synthetic}),
an additionally L1$_{\mathrm{topo}}$ was defined: full network routing
(nodal balance as in L3) but with zero pipe losses. This isolates
the pure routing effect from the thermal-loss effect within the
L1$\to$L3 gap.

\paragraph{Demand concentration metrics.}
The topology effect depends on how unevenly demand is distributed
across nodes. The study quantifies this via (i)~the \textbf{Gini
coefficient} of nodal annual demand shares and (ii)~the
\textbf{top-3 demand share}. For the primary
case, Gini\,$= 0.41$ and top-3 share\,$= 51.7\,\%$, indicating
moderate spatial concentration. The synthetic parametric study uses a
normalised variance index HI as a discrete parameter
(Low/Med/High); its definition and mapping to Gini are given in
\Cref{app:hi_definition}.

%% ============================================================
%% 3.2 BASIC MILP
%% ============================================================
\subsection{Basic MILP formulation (L1/L2/L3)}
\label{sec:milp}

All five variants share the objective function and component models
below; L1/L2/L3 differ only in the energy-balance topology and
loss treatment.

\paragraph{Objective function.}
The optimiser minimises total annual operational cost, which
consists of electricity procurement and feed-in revenues, fuel
expenditure, heat-curtailment penalties, carbon charges, and a
demand charge on peak grid import:
%
\begin{equation}
  \min \; Z = C^{\mathrm{en}} + C^{\mathrm{fuel}}
    + C^{\mathrm{dump}} + C^{\ce{CO2}} + C^{\mathrm{dem}}
    + C^{\mathrm{pump}}
  \label{eq:objective}
\end{equation}
%
The pumping cost $C^{\mathrm{pump}} = 0$ for L1/L2/L3 and enters
only in $\Lplus$/$\Lnl$ (\Cref{sec:extended}). The individual
terms are defined as follows, where $\Delta t = 1\,\mathrm{h}$ is
the time-step length and $t$ indexes all 8{,}760 hours of the year:
%
\begin{align*}
C^{\mathrm{en}}   &= \Delta t \sum_t
  \bigl(P_t^{\mathrm{buy}} \lambda_t^{\mathrm{buy}}
        - P_t^{\mathrm{sell}} \lambda_t^{\mathrm{sell}}\bigr) \\
C^{\mathrm{fuel}} &= \Delta t \sum_{g,t}
  F_{g,t}\, \lambda_g^{\mathrm{fuel}} \\
C^{\mathrm{dump}} &= \Delta t \sum_{n,t}
  Q_{n,t}^{\mathrm{dump}}\, \lambda^{\mathrm{dump}} \\
C^{\ce{CO2}}      &= \frac{\lambda^{\ce{CO2}}}{1000} \sum_t
  \Bigl(E_t^{\ce{CO2},\mathrm{grid}} +
        \sum_g E_{g,t}^{\ce{CO2},\mathrm{fuel}}\Bigr) \\
C^{\mathrm{dem}}  &= \lambda^{\mathrm{dem}} f^{\mathrm{yr}}
  P^{\mathrm{peak}}
\end{align*}
%
Here, $P_t^{\mathrm{buy}}$ and $P_t^{\mathrm{sell}}$
[MW] are the grid import and export at time~$t$;
$\lambda_t^{\mathrm{buy}}$ and $\lambda_t^{\mathrm{sell}}$
[EUR/MWh] are the corresponding prices; $F_{g,t}$ [MW]
is the fuel input to generator~$g$ at heating value basis;
$\lambda_g^{\mathrm{fuel}}$ [EUR/MWh] is its specific fuel cost;
$Q_{n,t}^{\mathrm{dump}}$ [MW] is curtailed heat at node~$n$,
penalised at $\lambda^{\mathrm{dump}}$ [EUR/MWh];
$\lambda^{\ce{CO2}}$ [EUR/t$_{\mathrm{CO_2}}$] is the carbon
price; $f^{\mathrm{yr}}$ a fractional annualisation factor; and
$P^{\mathrm{peak}}$ [MW] the annual peak grid import entering the
demand charge $\lambda^{\mathrm{dem}}$ [MWh$^{-1}$].

\paragraph{Energy balance.}
L1 enforces a single bus-level balance across all generators and
storage units, ignoring network losses and spatial routing:
%
\begin{equation}
  \sum_{g} Q_{g,t}^{\mathrm{th}}
    + \sum_{s} Q_{s,t}^{\mathrm{dis}}
  = Q_t^{\mathrm{dem}} + Q_t^{\mathrm{dump}}
    + \sum_{s} Q_{s,t}^{\mathrm{ch}}
  \;\; \forall\, t
  \label{eq:balance_L1}
\end{equation}
%
where $Q_{g,t}^{\mathrm{th}}$ [MW] is the thermal output of
generator~$g$, $Q_{s,t}^{\mathrm{dis}}$ and $Q_{s,t}^{\mathrm{ch}}$
[MW] are the discharge and charge of storage unit~$s$, and
$Q_t^{\mathrm{dem}}$ [MW] is the total network heat demand.

L2/L3/$\Lplus$/$\Lnl$ instead enforce a \emph{nodal} balance at
each junction~$n$, accounting for thermal losses along incoming
pipes and the spatial distribution of assets and demands:
%
\begin{align}
  &\sum_{p \in \mathcal{P}_n^{\mathrm{in}}}
    \bigl(Q_{p,t} - Q_{p,t}^{\mathrm{loss,pipe}}\bigr)
  + \sum_{g \in \mathcal{G}_n} Q_{g,t}^{\mathrm{th}}
  + \sum_{s \in \mathcal{S}_n} Q_{s,t}^{\mathrm{dis}} \notag \\
  &= Q_{n,t}^{\mathrm{dem}} + Q_{n,t}^{\mathrm{dump}}
  + \sum_{s \in \mathcal{S}_n} Q_{s,t}^{\mathrm{ch}}
  + \sum_{p \in \mathcal{P}_n^{\mathrm{out}}} Q_{p,t}
  \;\; \forall\, n,\, t
  \label{eq:balance_nodal}
\end{align}
%
Here $\mathcal{P}_n^{\mathrm{in}}$ and $\mathcal{P}_n^{\mathrm{out}}$
are the sets of pipes entering and leaving node~$n$;
$Q_{p,t}$ [MW] is the heat flow in pipe~$p$; and
$Q_{p,t}^{\mathrm{loss,pipe}}$ [MW] is the corresponding thermal
loss (defined below). The contrast between
\cref{eq:balance_L1,eq:balance_nodal} encodes the central
experimental variable: the L1 model's missing pipe-loss and
spatial-flow terms are the mechanism driving the L1\,$\to$\,L2
cost gap.

The global electricity balance and grid mutual-exclusivity
constraints close the system. All electrical generators, heat
pumps (power-to-heat units~$r$), and pumps must be balanced
against grid import, export, and in-network consumption:
%
\begin{align}
  P_t^{\mathrm{buy}}
    &+ \sum_{g \in \mathcal{G}^{\mathrm{CHP}}} P_{g,t}^{\mathrm{el}}
  = \sum_{h \in \mathcal{H}} P_{h,t}^{\mathrm{el}}
    + \sum_{r \in \mathcal{R}} P_{r,t}^{\mathrm{el}} \notag \\
    &+ P_t^{\mathrm{pump}} + P_t^{\mathrm{sell}}
  \quad \forall\, t
  \label{eq:balance_el} \\[4pt]
  P_t^{\mathrm{buy}} &\leq M^{\mathrm{grid}}\, z_t,
  \;\;
  P_t^{\mathrm{sell}} \leq M^{\mathrm{grid}}\, (1 - z_t)
  \;\; \forall\, t
  \label{eq:grid_bigm}
\end{align}
%
where $P_{g,t}^{\mathrm{el}}$ [MW] is the CHP electrical output,
$P_{h,t}^{\mathrm{el}}$ [MW] the heat pump electrical input,
$P_{r,t}^{\mathrm{el}}$ [MW], the electrode boiler input, and
$P_t^{\mathrm{pump}}$ [MW] the network pumping power.
The binary variable $z_t \in \{0,1\}$ enforces that the grid
operates in either import or export mode at any hour;
$M^{\mathrm{grid}}$ is a Big-M constant set to the maximum
grid capacity.

\paragraph{Thermal loss model (L2/L3).}
\label{sec:model_losses}
For each pipe~$p$, steady-state conductive heat losses are computed
separately for the supply and return circuit. Because the supply
temperature is determined by the heating curve rather than
optimised, the losses are \emph{fixed parameters}. This is a key
simplification that preserves MILP linearity:
%
\begin{align}
  Q_{p,t}^{\mathrm{loss,sup}} &= \frac{U_p\, L_p\,
    \bigl(T_{\mathrm{sup}}^{\mathrm{nom}}(t)
    - T_{\mathrm{ground}}(t)\bigr)}{10^6}
  \label{eq:loss_supply} \\[4pt]
  Q_{p,t}^{\mathrm{loss,ret}} &= \frac{U_p\, L_p\,
    \bigl(T_{\mathrm{ret}}^{\mathrm{nom}}(t)
    - T_{\mathrm{ground}}(t)\bigr)}{10^6}
  \label{eq:loss_return}
\end{align}
%
where $U_p$ [MW\,/(m\,K)] is the linear heat transfer
coefficient of pipe~$p$ (determined by pipe diameter and insulation
class), $L_p$ [m] its length, $T_{\mathrm{sup}}^{\mathrm{nom}}(t)$
and $T_{\mathrm{ret}}^{\mathrm{nom}}(t)$ [\si{\celsius}] the nominal
supply and return temperatures prescribed by the heating curve, and
$T_{\mathrm{ground}}(t)$ [\si{\celsius}] the ambient ground
temperature. The factor $10^6$ converts from W to MW. For L1,
all pipe losses are identically zero.

\paragraph{Heat pump model.}
Each heat pump~$h$ can recover industrial waste heat (WHR source,
temperature-varying) or fall back on a default low-temperature
source (constant COP). The total thermal output is the sum of both
contributions:
%
\begin{align}
  Q_{h,t} &= Q_{h,t}^{\mathrm{WHR}} + Q_{h,t}^{\mathrm{def}}
  \label{eq:hp_split} \\[4pt]
  P_{h,t}^{\mathrm{el}} &= \frac{Q_{h,t}^{\mathrm{WHR}}}{
    \mathrm{COP}_{h,t}^{\mathrm{WHR}}}
    + \frac{Q_{h,t}^{\mathrm{def}}}{\mathrm{COP}_h^{\mathrm{def}}}
  \label{eq:hp_pel}
\end{align}
%
\Cref{eq:hp_split} splits the thermal output $Q_{h,t}$ [MW]
between the two source circuits; \Cref{eq:hp_pel} computes the
corresponding electrical demand $P_{h,t}^{\mathrm{el}}$ [MW] using
the time-varying waste-heat $\mathrm{COP}_{h,t}^{\mathrm{WHR}}$
and the constant default $\mathrm{COP}_h^{\mathrm{def}}$.
The waste-heat COP is pre-computed offline from an analytical
Lorenz-cycle model (\Cref{app:cop}); this avoids introducing
bilinear terms in the optimisation  model.

Operation is bounded by a minimum load fraction
$\phi_h^{\mathrm{min}}$ and the installed capacity $\bar{Q}_h$,
controlled by the on/off binary $u_{h,t} \in \{0,1\}$:
%
\begin{equation}
  \phi_h^{\mathrm{min}}\, \bar{Q}_h\, u_{h,t}
  \leq Q_{h,t} \leq \bar{Q}_h\, u_{h,t}
  \quad \forall\, h,\; \forall\, t
  \label{eq:hp_cap}
\end{equation}
%
When $u_{h,t} = 0$, the heat pump is off, and both outputs are
forced to zero; when $u_{h,t} = 1$ the output must reach the
minimum load fraction $\phi_h^{\mathrm{min}}$, preventing
uneconomic part-load cycling.

\paragraph{Thermal storage.}
The state of charge~$E_{s,t}$ [MWh] of storage unit~$s$ evolves
according to a discrete-time energy balance that accounts for
self-discharge, charging efficiency $\eta_s^{\mathrm{ch}}$, and
discharging efficiency $\eta_s^{\mathrm{dis}}$:
%
\begin{equation}
  E_{s,t} = (1 - \alpha_s)^{\Delta t}\, E_{s,t-1}
    + \eta_s^{\mathrm{ch}}\, Q_{s,t}^{\mathrm{ch}}\, \Delta t
    - \frac{Q_{s,t}^{\mathrm{dis}}\, \Delta t}{\eta_s^{\mathrm{dis}}}
  \label{eq:storage_soc}
\end{equation}
%
The self-discharge coefficient $\alpha_s$
[$h^{-1}$] models standing heat losses of the storage vessel;
the exponential form $(1-\alpha_s)^{\Delta t}$ provides a
consistent discretisation independent of the time step. A cyclic
boundary condition $E_{s,|\mathcal{T}|} = E_{s,0}$ enforces
energy-neutral operation over the full simulation year, preventing
artificial depletion or accumulation.

\paragraph{Remaining generation assets and emissions.}
Gas boiler, biomass boiler, electrode boiler, and CHP follow
standard linear formulations (capacity-bounded, constant
efficiency; detailed in \Cref{app:component_models}).

Grid and fuel emissions are computed as the product of energy
consumed and the respective emission intensity. This allows the
optimiser to trade off electricity procurement against fuel
combustion in a time-varying carbon context:
%
\begin{equation}
  E_t^{\ce{CO2},\mathrm{grid}} =
    P_t^{\mathrm{buy}}\, e_t^{\mathrm{grid}}\, \Delta t,
  \qquad
  E_{g,t}^{\ce{CO2},\mathrm{fuel}} =
    F_{g,t}\, e_g^{\mathrm{fuel}}\, \Delta t
  \label{eq:emissions}
\end{equation}
%
where $e_t^{\mathrm{grid}}$ [kg\,MWh$^{-1}$] is the hourly
marginal grid emission factor and $e_g^{\mathrm{fuel}}$
[kg\,MWh$^{-1}$] the fuel-specific emission intensity.
Both values are fixed input parameters; the optimiser is thus
carbon-aware but does not optimise the grid emission factor.

%% ============================================================
%% 3.3 EXTENDED PHYSICS (L+ and Lnl)
%% ============================================================
\subsection{Extended physics formulation ($\Lplus$ / $\Lnl$)}
\label{sec:extended}

$\Lplus$ and $\Lnl$ augment L3 with three additional physical
phenomena: pressure-dependent pumping losses, temperature-dependent
heat propagation along pipes, and transport delay. They share
identical physics scope but differ in mathematical treatment:
$\Lplus$ linearises all nonlinear terms for MILP solvability, while
$\Lnl$ retains native quadratic and bilinear constraints (MIQCP).
All PWL approximations in $\Lplus$ use SOS2 machinery 
with $K{=}3$ segments and $K{+}1=4$ support points. A complete side-by-side
comparison is provided in \Cref{app:extended_formulations}.

\paragraph{Pressure drop and pumping power.}
\label{sec:model_pressure}
Fluid friction in pipe~$p$ is governed by the Darcy--Weisbach
equation \cite{vanderHeijde2017}, which relates pressure drop $\Delta p_p$ [Pa] to the
square of the mass flow rate $\dot{m}_p$ [kg\,s$^{-1}$]:
%
\begin{equation}
  \Delta p_{p} = f_D\, \frac{L_p}{D_p}\,
    \frac{8\,\dot{m}_p^2}{\rho \pi^2 D_p^4}
  \label{eq:darcy}
\end{equation}
%
where $f_D$ [-] is the Darcy friction factor, $D_p$ [m] the inner
pipe diameter, and $\rho$ [kg\,m$^{-3}$] the fluid density.
Since supply and return temperatures are prescribed by the heating
curve, the mass flow rate is a linear function of the pipe heat
flow $Q_{p,t}$ [MW]:
%
\begin{equation}
  \dot{m}_{p,t} = \frac{Q_{p,t}}{
    c_p \bigl(T_{\mathrm{sup}}^{\mathrm{nom}}(t)
    - T_{\mathrm{ret}}^{\mathrm{nom}}(t)\bigr)}
  \label{eq:massflow}
\end{equation}
%
where $c_p$ [J\,(kg\,K)$^{-1}$] is the specific heat capacity of
water. Substituting \Cref{eq:massflow} into \Cref{eq:darcy} reveals
that pressure drop is \emph{quadratic} in $Q_{p,t}$; it can be
written compactly as $\Delta p_{p,t} = R_p \cdot Q_{p,t}^2$, where
$R_p$ [Pa\,MW$^{-2}$] is the pipe hydraulic resistance.

In $\Lnl$ this quadratic relation enters directly as a native
quadratic constraint. In $\Lplus$, it is approximated by a
piecewise-linear (PWL) function over $K$ equal segments using SOS2
weights $w_{p,t,k} \in [0,1]$:
%
\begin{equation}
  \Delta p_{p,t}^{\mathrm{PWL}} = \sum_{k=1}^{K}
    m_k^{(p)} \cdot w_{p,t,k}
  \label{eq:pressure_pwl}
\end{equation}
%
where $m_k^{(p)}$ are the slopes on each segment. The maximum
approximation error of this PWL is (derivation in \Cref{app:pwl}):
%
\begin{equation*}
  \varepsilon_{\Delta p}^{\max}
    = \frac{R_p\, \bar{Q}_p^2}{4K^2}
    \approx 2.8\,\% \text{ of peak pressure drop for } K{=}3
\end{equation*}
%
where $\bar{Q}_p$ [MW] is the rated pipe capacity.

The network pumping power $P_t^{\mathrm{pump}}$ [MW] is the sum
of hydraulic work across all pipes, normalised by the pump
efficiency $\eta_{\mathrm{pump}}$:
%
\begin{equation}
  P_t^{\mathrm{pump}} = \frac{1}{\eta_{\mathrm{pump}}}
    \sum_{p \in \mathcal{P}}
    \frac{\dot{m}_{p,t} \cdot \Delta p_{p,t}}{\rho}
  \label{eq:pump_power}
\end{equation}
%
Because this involves $\dot{m}_{p,t} \propto Q_{p,t}$ and
$\Delta p_{p,t} \propto Q_{p,t}^2$, pumping power is cubic in
$Q_{p,t}$. In $\Lplus$, the cubic relation is directly
piecewise-linearised. In $\Lnl$, it is modelled via the bilinear
auxiliary variable $W_{p,t} = Q_{p,t}^2$, entering the objective
as $Q_{p,t} \cdot W_{p,t}$.

\paragraph{Temperature-dependent loss propagation.}
\label{sec:model_temp_prop}
In L3, pipe losses are computed from fixed nominal temperatures. In
$\Lplus$ and $\Lnl$, the actual outlet temperature is tracked as
the fluid travels along pipe~$p$. Under steady-state conditions,
the outlet temperature $T_{p,t}^{\mathrm{out}}$ [\si{\celsius}]
decays exponentially from the inlet temperature
$T_{n_1,t}^{\mathrm{sup}}$ [\si{\celsius}] towards the ground
temperature $T_{\mathrm{gr}}(t)$ [\si{\celsius}]:
%
\begin{equation}
  T_{p,t}^{\mathrm{out}} = T_{\mathrm{gr}}(t)
    + \bigl(T_{n_1,t}^{\mathrm{sup}} - T_{\mathrm{gr}}(t)\bigr)
    \cdot \underbrace{\exp\!\left(
      \frac{-U_p L_p}{\dot{m}_{p,t}\, c_p}
    \right)}_{\displaystyle\phi_p(\dot{m}_{p,t})}
  \label{eq:temp_prop}
\end{equation}
%
The decay factor $\phi_p \in (0,\,1]$ quantifies what fraction
of the inlet temperature excess is preserved at the outlet; it
approaches unity for high flow rates (little cooling) and zero
for near-zero flow (complete equalisation to ground temperature).
Two nonlinearities arise from this expression: the exponential
$\phi_p(\dot{m}_{p,t})$ and the bilinear product
$T_{n_1,t}^{\mathrm{sup}} \cdot \phi_p$.

Both $\Lplus$ and $\Lnl$ approximate $\phi_p$ with a shared
PWL over $K$ breakpoints in $\dot{m}$ (this shared error is
excluded from the measured linearisation gap). The bilinear
product is handled differently in each level. In $\Lnl$,
it enters directly as $T_{n_1,t}^{\mathrm{sup}} \cdot \phi_{p,t}^{\mathrm{PWL}}$.
In $\Lplus$, a first-order Taylor expansion around the nominal
operating point $(T^{\mathrm{nom}}(t),\, \phi_p^{\mathrm{nom}})$
linearises the product:
%
\begin{align}
  T_{p,t}^{\mathrm{out}} &\approx T_{\mathrm{gr}}(t)
    + \bigl(T^{\mathrm{nom}}(t) - T_{\mathrm{gr}}(t)\bigr)
      \cdot \phi_{p,t}^{\mathrm{PWL}} \notag \\
    &+ \phi_p^{\mathrm{nom}}
      \cdot \bigl(T_{n_1,t}^{\mathrm{sup}} - T^{\mathrm{nom}}(t)\bigr)
  \label{eq:temp_prop_lin}
\end{align}
%
The neglected second-order cross-term is bounded by
$|\Delta T_{\max}| \cdot |\Delta\phi_{\max}| < 2.5\,\text{K}$
(derivation in \Cref{app:taylor_derivation}).

\paragraph{Transport delay.}
\label{sec:model_delay}
At finite flow velocity, heat injected at the source at time~$t$
does not arrive at the consumer until later. The nominal
travel time $\tau_p$ [h] through pipe~$p$ is determined by the
pipe cross-sectional area $A_p = \pi D_p^2/4$ [m$^2$] and
the nominal mass flow $\dot{m}_p^{\mathrm{nom}}$ [kg\,s$^{-1}$]:
%
\begin{equation}
  \tau_p = \frac{L_p\, \rho\, A_p}{\dot{m}_p^{\mathrm{nom}}}
  \label{eq:delay_time}
\end{equation}
%
Discretising to the nearest integer number of time steps
$k_p = \mathrm{round}(\tau_p / \Delta t)$, the heat delivered
at the downstream end is the heat injected $k_p$ steps earlier,
reduced by the corresponding pipe losses:
%
\begin{equation}
  Q_{p,t}^{\mathrm{delivered}} =
    Q_{p,\,t-k_p} - Q_{p,\,t-k_p}^{\mathrm{loss,pipe}}
  \quad \forall\, p,\; t \geq k_p{+}1
  \label{eq:delayed_balance}
\end{equation}
%
Transport delay is active only in $\Lnl$. Under $\Lplus$ the
linearised pipe model bypasses the lag because the Taylor
linearisation of $\phi_p$ is performed at fixed nominal flow,
which implicitly assumes instantaneous transport
(cf.\ \Cref{tab:model_scope}).

\Cref{tab:lplus_vs_lnl} in \Cref{app:extended_formulations}
provides a complete side-by-side comparison of all extended-physics
treatments in $\Lplus$ and $\Lnl$.

%% ============================================================
%% 3.4 VALIDATION
%% ============================================================
\subsection{Validation protocol}
\label{sec:val_protocol}

The network underwent a technology upgrade (HP, electrode boiler,
TES) \emph{after} the measurement record was collected. Following
\citet{Kus2025} and \citet{Maldonado2024}, this work adopts a two-stage
strategy.

\paragraph{Stage~1: Direct network validation.}
The pipe model is validated against pre-upgrade monitoring data
(new assets set to zero capacity). Calibration follows the
branch-sequential $U_p$-tuning procedure of
\citet{Maldonado2024}. \Cref{tab:val_targets} lists the KPIs with
gates defined \emph{ex ante} from published practice. The
$T_{\mathrm{sup}}$ gates are assigned to $\Lnl$ because only this
variant propagates supply temperature along pipes; L3 uses fixed
nominal temperatures and therefore cannot predict downstream
temperature deviations.

\begin{table}[ht]
\centering
\begin{threeparttable}
\caption{Stage~1 validation KPIs and gates (defined ex ante).}
\label{tab:val_targets}
\footnotesize
\setlength{\tabcolsep}{2pt}
\begin{tabular*}{\columnwidth}{@{\extracolsep{\fill}} l l l l @{}}
\toprule
KPI & Model & Gate & Source \\
\midrule
$T_{\mathrm{sup}}$, mean
  & $\Lnl$ & ${<}\SI{1.5}{\kelvin}$
  & \citet{Kus2025} \\
$T_{\mathrm{sup}}$, worst
  & $\Lnl$ & ${<}\SI{2.5}{\kelvin}$
  & \citet{Maldonado2024} \\
$T_{\mathrm{ret}}$ at source
  & L3 & ${<}\SI{2.5}{\kelvin}$
  & \citet{Kus2025} \\
Flow (MAPE)
  & L3 & n/a\tnote{a} & -- \\
Energy balance
  & L3 & ${<}2\,\%$ & Conserv. \\
\bottomrule
\end{tabular*}
\begin{tablenotes}\footnotesize
  \item[a] ${\pm}15\,\%$ instrument error; no gate.
\end{tablenotes}
\end{threeparttable}
\end{table}

\paragraph{Stage~2: Necessary-condition verification.}
Dispatch of new assets is verified through plausibility bounds:
HP COP within $[2.5,\,5.5]$; electrode boiler price-response
correlation $r > 0.3$; per-step energy closure ${<}\,1\,\%$; TES
SOC within $[0.05,\,0.95] \cdot \bar{E}_s$. These are necessary,
not sufficient; post-upgrade data are recommended for follow-up.

\paragraph{Structural sanity checks.}
Two additional checks ensure model consistency: (i)~$\Lnl$ with
all quadratic constraints deactivated reproduces L3's objective
within 0.01\,\%; (ii)~the relaxation property
$\mathrm{cost}(\mathrm{L1}) \leq \mathrm{cost}(\mathrm{L3})$ holds
for every synthetic instance.

%% ============================================================
%% 3.5 IMPLEMENTATION
%% ============================================================
\subsection{Implementation}
\label{sec:model_impl}

All five variants are implemented in Python using PYOMO 
\cite{pyomo}, instantiated from a single parameterised codebase
via YAML configuration files. Solver: Gurobi~13.0 with MIP gap
${\leq}\,0.5\,\%$ for both MILP and MIQCP
(\texttt{NonConvex=2}); identical random seeds. $\Lnl$ is
warm-started from the $\Lplus$ solution.

The $\Lplus$ formulation adds ${\approx}\,753{,}000$ variables
relative to L3, of which ${\approx}\,368{,}000$ are binary
(SOS2 mode selection for piecewise linearisation) and
${\approx}\,385{,}000$ continuous; $\Lnl$ adds
${\approx}\,245{,}000$ quadratic constraints. A single
744\,h MIQCP solve requires ${>}\,4$\,h at 1.99\,\% gap; the
full-year formulation is intractable on single-workstation hardware from the 
University of Stuttgart. The synthetic parametric study is therefore restricted
to L1/L2/L3/$\Lplus$, while the linearisation error is
characterised on the primary case.

%% ============================================================
%% 4. CASE STUDIES
%% ============================================================
\section{Case study}
\label{sec:case_studies}
The case study comprises data from a real industrial heating network and several synthetic networks to generate generalisable results.

\subsection{Primary case: industrial district heating network}
\label{sec:case_stadtbach}

The primary case is an industrial DH network operated by a
medium-sized utility in southern Germany. The network has a radial
tree structure with a central production node~(j$_1$) and two
distribution arms (15~junctions, 14~pipe segments, 27~consumers;
\Cref{fig:topology}). All generation assets (CHP, gas boiler, biomass boiler, heat pump, and electrode boiler) are modelled at the central production node j1. The heat pump recovers industrial waste heat via a time-varying Lorenz-cycle COP (Appendix A.1), independent of its position in the network graph.

Supply temperature follows a piecewise-linear heating curve from
\SI{70}{\celsius} (low load) to \SI{100}{\celsius} (peak load),
with corresponding return temperatures of \SI{55}{\celsius} to
\SI{40}{\celsius} (\Cref{app:heating_curve}). Simulated annual
thermal pipe losses are approximately 5\,\% of network heat
demand. Absolute demand values are anonymised per non-disclosure
agreement, and the key results are reported as relative differences.
\Cref{tab:case_summary} consolidates the key parameters.

L2 aggregates the 15~L3 junctions into seven~zones by merging
adjacent nodes sharing a branch segment
(\Cref{app:l2_zone_mapping}). The total $\sum U_p L_p$ is
preserved exactly (\SI{1,011}{\watt\per\kelvin}).

\begin{figure}[ht]
  \centering
  \includegraphics[width=\linewidth]{paper1_dh_fidelity/F2_network_topology.png}
  \caption{Network topology. Assets at j$_1$ (CHP, TES, boilers, HP, electrode boiler). Transport delays annotated.}
  \label{fig:topology}
\end{figure}

\begin{table}[ht]
\centering
\caption{Primary case study: consolidated parameters.}
\label{tab:case_summary}
\begin{threeparttable}
\scriptsize
\setlength{\tabcolsep}{2pt}
\begin{tabular*}{\columnwidth}{@{\extracolsep{\fill}} l l l @{}}
\toprule
Parameter & Value & Notes \\
\midrule
\multicolumn{3}{@{}l}{\textit{Network}} \\
Nodes/consumers/pipes
  & 15/27/14 & Radial tree \\
Total pipe length
  & \SI{3255}{\metre} & Longest: \SI{2125}{\metre} \\
Max.\ $\tau$
  & ${\approx}\SI{35}{\minute}$ & $k_p{=}1$ \\
\midrule
\multicolumn{3}{@{}l}{\textit{Thermal}} \\
$T_{\mathrm{sup}}$/$T_{\mathrm{ret}}$
  & 70--100/40--\SI{55}{\celsius}
  & PWL heating curve \\
$T_{\mathrm{ground}}$
  & \SI{10}{\celsius} & Fixed \\
$U$-value (sup/ret)
  & 0.28--0.34\,\si{\watt\per\metre\per\kelvin}
  & Per-pipe \\
\midrule
\multicolumn{3}{@{}l}{\textit{Generation at~j$_1$}} \\
Gas boiler/biomass
  & 13/\SI{3.3}{\mega\watt_{\mathrm{th}}}
  & $\eta{=}0.90$/$0.85$ \\
CHP (th/el)
  & 0.2/\SI{0.1}{\mega\watt}
  & Backpressure\tnote{a} \\
Thermal storage
  & \SI{500}{\mega\watt\hour}/\SI{30}{\mega\watt}
  & Cyclic BC \\
HP (WHR)/electrode
  & 5/\SI{5}{\mega\watt_{\mathrm{th}}}
  & Lorenz $\eta{=}0.6$/$0.95$ \\
\midrule
\multicolumn{3}{@{}l}{\textit{Demand}} \\
Peak/Gini/top-3
  & ${\approx}\SI{12}{\mega\watt}$/0.41/51.7\,\%
  & Moderate concentration \\
\bottomrule
\end{tabular*}
\begin{tablenotes}\scriptsize
  \item[a] CHP supplies ${<}2\,\%$ of annual output.
\end{tablenotes}
\end{threeparttable}
\end{table}

\subsection{Synthetic parametric network}
\label{sec:case_synthetic}

To generalise beyond the primary case, a synthetic tree network is
parameterised by four properties (\Cref{tab:synthetic}). Of the $3^4 = 81$ parameter combinations, 36 are retained  after feasibility filtering. Configurations are classified 
as infeasible when the LP relaxation of the L1 copperplate 
model has no solution, which occurs primarily at extreme 
parameter combinations: high demand heterogeneity 
(HI\,$=\,0.8$) combined with long pipe networks 
($L\,{=}\,15$\,km) and low storage capacity 
($\tau_s\,{=}\,2$\,h) creates supply demand imbalances 
that no asset configuration can resolve. All 36 retained 
configurations are feasible across all five model levels 
(L1 through $\Lplus$).

\begin{table}[ht]
\centering
\caption{Synthetic network parameter space (36~feasible).}
\label{tab:synthetic}
\footnotesize
\begin{tabular*}{\columnwidth}{@{\extracolsep{\fill}} l c c c l @{}}
\toprule
Parameter & Low & Med & High & Rationale \\
\midrule
Pipe [km]    & 1   & 5   & 15  & Campus--urban \\
Demand HI    & 0.1 & 0.4 & 0.8 & Uniform--conc. \\
Node count   & 5   & 15  & 30  & Sparse--dense \\
Storage [h]  & 2   & 6   & 12  & Buffer--shift \\
\bottomrule
\end{tabular*}
\end{table}

The parametric study thus comprises $36 \times 5$ instances
(L1/L1$_{\mathrm{topo}}$/L2/L3/$\Lplus$ per configuration)
plus 4~primary-case levels, totalling 184~optimisation  runs.

\paragraph{Additive gap decomposition.}
The parametric study includes a zero-physics topology stage
L1$_{\mathrm{topo}}$ (full network routing, zero losses) to
decompose the cumulative L1$\to$L3 gap:
%
\begin{equation}
\label{eq:cost_decomposition}
\begin{split}
  \underbrace{\mathrm{cost}(L3) - \mathrm{cost}(L1)}_{\text{total gap}}
  \;=\;&
  \underbrace{\mathrm{cost}(L1_{\mathrm{topo}}) - \mathrm{cost}(L1)}_{\text{routing}} \\[6pt]
  &+\;
  \underbrace{\mathrm{cost}(L2) - \mathrm{cost}(L1_{\mathrm{topo}})}_{\text{heat loss}} \\[6pt]
  &+\;
  \underbrace{\mathrm{cost}(L3) - \mathrm{cost}(L2)}_{\text{pressure drop}}
\end{split}
\end{equation}
%
To enable this three-way decomposition, the synthetic study
redefines the physics scope of L2 and L3 relative to the
primary-case taxonomy (\Cref{tab:model_scope}): synthetic~L2
includes steady-state thermal losses but no pressure drop,
while synthetic~L3 adds Darcy--Weisbach pressure drop on top.
\Cref{tab:synth_level_mapping} clarifies this mapping.
The L3$\to\Lplus$ step in the synthetic study therefore
isolates only transport-delay activation (which produces zero
cost change at hourly resolution when
$\tau_{\max} < 1$\,h); the pressure-drop effect itself is
captured within the L2$\to$L3 decomposition term. For the
primary case, pressure drop is excluded from L3 and enters
exclusively at $\Lplus$, consistent with
\Cref{tab:model_scope}.

\begin{table}[ht]
\centering
\caption{Physics-scope mapping between primary case and
synthetic parametric study. The reassignment enables the
three-term additive decomposition (\Cref{eq:cost_decomposition}).}
\label{tab:synth_level_mapping}
\begin{threeparttable}
\footnotesize
\begin{tabular*}{\columnwidth}{@{\extracolsep{\fill}}
  >{\raggedright\arraybackslash}p{0.14\columnwidth}
  >{\raggedright\arraybackslash}p{0.34\columnwidth}
  >{\raggedright\arraybackslash}p{0.14\columnwidth}
  >{\raggedright\arraybackslash}p{0.20\columnwidth} @{}}
\toprule
Synth.\ label & Physics content & Primary equiv.
  & Decomp.\ role \\
\midrule
L1
  & No losses, no $\Delta p$
  & L1
  & Baseline \\
L1$_{\mathrm{topo}}$
  & Routing only, no losses
  & --
  & Routing ref. \\
L2
  & Thermal losses, no $\Delta p$
  & $\approx$\,L3\tnote{a}
  & Heat-loss ref. \\
L3
  & Thermal losses + $\Delta p$
  & $\approx$\,$\Lplus$\tnote{b}
  & $\Delta p$ target \\
$\Lplus$
  & Losses + $\Delta p$ + delay
  & --
  & Delay check \\
\bottomrule
\end{tabular*}
\begin{tablenotes}\scriptsize
  \item[a] Equivalent in loss treatment; differs in node
    aggregation (synthetic L2 uses full node resolution, not
    the 7-zone aggregation of the primary-case L2).
  \item[b] Without Taylor-linearised temperature propagation;
    pressure drop enters as PWL (Darcy--Weisbach) in both.
\end{tablenotes}
\end{threeparttable}
\end{table}


%% ============================================================
%% 5. RESULTS
%% ============================================================
\section{Results}
\label{sec:results}
In the results section, the validation results are presented first, followed by the topology and thermo-hydraulic influencing factors. Finally, the generalisability, sensitivity, and computational performance are examined. 


%% ============================================================
%% 5.1 VALIDATION RESULTS
%% ============================================================
\subsection{Validation results}
\label{sec:results_validation}

\Cref{tab:validation_stage1} reports the Stage~1 results; all gates
defined ex ante (\Cref{tab:val_targets}) are met.
\Cref{fig:val_timeseries} shows a representative winter week
(Dec.~22--28) with dispatch composition by asset class (top panel)
and thermal storage state-of-charge (bottom panel); measured
boundary conditions are reproduced within the prescribed gates.
The elevated flow MAPE (36\,\%) is attributable to
pressure-differential metering at partial load; the 1.23\,\%
annual energy error confirms that energy-weighted flow is correctly
reproduced.

\begin{figure}[!htbp]
  \centering
  \includegraphics[width=\linewidth]{paper1_dh_fidelity/FV1_validation_timeseries.png}
  \caption{Stage~1 validation, representative winter week
  (Dec.~22--28): dispatch composition (top) and TES
  state-of-charge (bottom).}
  \label{fig:val_timeseries}
\end{figure}

Node j$_1$ is the production source and serves as boundary
condition (measured inlet temperature); it is excluded from the
temperature-error evaluation. Of the 14~remaining network nodes,
6 enter the gate evaluation (``validation nodes'',
cf.\ \Cref{tab:val_targets}); the remaining 8 are excluded for two
distinct reasons (per-node MAE breakdown in
\Cref{app:per_node_mae}). Four nodes (j$_6$, j$_7$, j$_8$,
j$_{14}$) served as \emph{calibration nodes} during the
branch-sequential $U_p$-tuning procedure and are therefore
in-sample residuals; their group mean is \SI{2.32}{\celsius}, and
node j$_7$ at \SI{3.07}{\celsius} would exceed the worst-node gate
of \SI{2.5}{\celsius} if held out for validation. This in-sample
deterioration is expected and motivates the strict
calibration/validation split adopted here. The remaining four nodes
(j$_2$, j$_3$, j$_4$, j$_5$) are excluded for measurement reasons
(mixing-valve offset above \SI{2}{\celsius} or absence of a single
junction reference among multiple downstream consumers); two of
them (j$_3$: \SI{0.58}{\celsius}; j$_4$: \SI{0.37}{\celsius}) are
in fact the two best-performing nodes in the network and would pass
either gate comfortably. Mean MAE across all 14~nodes is
\SI{1.68}{\celsius}.

\begin{table}[!htbp]
\centering
\caption{Stage~1 validation results.}
\label{tab:validation_stage1}
\begin{threeparttable}
\footnotesize
\begin{tabular*}{\columnwidth}{@{\extracolsep{\fill}} l c c c @{}}
\toprule
KPI & Model & Result & Gate \\
\midrule
$T_{\mathrm{sup}}$, mean\tnote{a}
  & $\Lnl$ & \SI{1.32}{\celsius} & ${<}\SI{1.5}{\celsius}$ \\
$T_{\mathrm{sup}}$, worst\tnote{a}
  & $\Lnl$ & \SI{2.27}{\celsius} & ${<}\SI{2.5}{\celsius}$ \\
$T_{\mathrm{ret}}$ at source
  & L3     & \SI{2.08}{\celsius} & ${<}\SI{2.5}{\celsius}$ \\
Flow (MAPE)
  & L3     & $36\,\%$            & n/a \\
Annual energy
  & L3     & $1.23\,\%$          & ${<}2\,\%$ \\
\bottomrule
\end{tabular*}
\begin{tablenotes}\footnotesize
  \item[a] 744\,h winter window; 6~of 14~nodes used as
    validation nodes (4~calibration, 4~excluded for measurement
    reasons; see \Cref{app:per_node_mae}).
\end{tablenotes}
\end{threeparttable}
\end{table}

Stage~2 plausibility checks all pass: observed HP COP within
$[2.5,\,3.0]$ (bounded by the moderate waste-heat source
temperature of 25--\SI{35}{\celsius} and the high sink temperature
of 70--\SI{100}{\celsius}); electrode-boiler price-response
correlation $r = 0.47$; per-step energy closure below 0.3\,\%; TES
SOC in $[0.18,\,0.94] \cdot \bar{E}_s$. Solver consistency:
$\Lnl$ with quadratic constraints deactivated reproduces L3 within
$0.004\,\%$.

\paragraph{Independent forward-physics cross-check.}
As complementary validation, a boundary-condition method (BCM)
propagates the measured source temperature at j$_1$ through the
steady-state pipe-loss model without dispatch optimisation,
isolating pipe physics from the optimiser. At the far-end node
j$_{15}$ (path \SI{2125}{\metre}), a global trunk $U$-multiplier
(${\times}1.330$) is calibrated on the Oct--Feb winter subset
across two partial heating seasons ($n = 3173$\,h). The resulting
error metrics are MAE\,$=$\,\SI{1.56}{\celsius},
RMSE\,$=$\,\SI{2.15}{\celsius}, and a residual positive bias of
\SI{+0.72}{\celsius}, indicating mild systematic overestimation
consistent with unmodelled flow-rate variability. Sixty percent of
hourly predictions fall within ${\pm}$\SI{1.5}{\celsius}; the
remaining deviations correlate with rapid flow transients
(partial-load switching) that the steady-state pipe model does not
capture. The mean MAE (\SI{1.56}{\celsius}) nevertheless remains
within the Stage~1 gate. \Cref{fig:val_scatter} presents the
scatter plot of simulated versus measured supply temperature;
\Cref{fig:val_timeseries_jan} shows the corresponding January
time-series with the ${\pm}$\SI{1.5}{\celsius} confidence band.
The calibration approach differs methodologically from the $\Lnl$
MIQCP-output validation reported in \Cref{tab:validation_stage1};
details are given in \Cref{app:bcm_calibration}.

\begin{figure}[!htbp]
  \centering
  \includegraphics[width=\linewidth]{paper1_dh_fidelity/F_validation_scatter_tsup.png}
  \caption{BCM cross-check: simulated vs.\ measured
  $T_{\mathrm{sup}}$ at j$_{15}$. Dotted lines mark the
  ${\pm}$\SI{1.5}{\celsius} gate.}
  \label{fig:val_scatter}
\end{figure}

\begin{figure}[!htbp]
  \centering
  \includegraphics[width=\linewidth]{paper1_dh_fidelity/F_validation_timeseries_2026_01.png}
  \caption{BCM cross-check (January): $T_{\mathrm{sup}}$
  time-series at j$_{15}$. The shaded band marks the
  ${\pm}$\SI{1.5}{\celsius} gate. Underneath is the ambient
  temperature for the week.}
  \label{fig:val_timeseries_jan}
\end{figure}

%% ============================================================
%% 5.2 TOPOLOGY EFFECT
%% ============================================================
\subsection{Topology effect: L1 vs.\ L2 vs.\ L3 (RQ1)}
\label{sec:results_cost}

\Cref{tab:cost_co2} reports the topology comparison under identical
basic physics. The cumulative L1$\to$L3 gap reaches 13.0\,\% in
operational cost and 74\,t\,\ce{CO2}/year. The L1$\to$L2 step
accounts for 10.4\,\% (63\,t\,\ce{CO2}/yr), while
L2$\to$L3 contributes the remaining 2.6\,\% (11\,t\,\ce{CO2}/yr).

The L1$\to$L2 gap arises from two mechanisms: (1)~L1 ignores
network losses entirely,  the additive decomposition ( \Cref{sec:case_synthetic}) confirms that thermal losses, not routing or locational siting, drive the cumulative gap.

\begin{table}[!htbp]
\centering
\caption{Topology comparison (basic physics, L3 = reference).}
\label{tab:cost_co2}
\begin{threeparttable}
\footnotesize
\begin{tabular*}{\columnwidth}{@{\extracolsep{\fill}} l c c c c c @{}}
\toprule
Level & $\Delta$Cost & $\Delta$Fuel &
  $\Delta$Elec. & $\Delta$\ce{CO2} & t/yr \\
\midrule
L1 & $-13.0\%$ & $-12.9\%$ & $-14.0\%$ & $-8.7\%$ & $772$ \\
L2 & $-2.6\%$  & $-2.4\%$  & $-3.2\%$  & $-1.3\%$ & $835$ \\
L3 & $(ref.)$  & $(ref.)$  & $(ref.)$  & $(ref.)$ & $846$ \\
\bottomrule
\end{tabular*}
\end{threeparttable}
\end{table}


%% ============================================================
%% 5.3 PHYSICS + LINEARISATION
%% ============================================================
\subsection{thermo-hydraulic effect and linearisation error (RQ2 + RQ3)}
\label{sec:results_extended}

\paragraph{Extended physics: L3\,$\to$\,$\Lplus$.}
Adding linearised pressure drop and temperature propagation
increases annual cost by $+0.11\,\%$, attributable entirely to
pumping power (+255\,EUR/yr on 2.4\,MWh additional pump
electricity). The thermo-hydraulic gap is thus two orders of
magnitude smaller than the topology gap. This answers RQ2 with a
clear negative: extended hydraulic and thermal physics do not
materially change dispatch when topology is already resolved at the
scale of the primary case.

\begin{table}[!htbp]
\centering
\caption{Thermo-hydraulics and linearisation comparison.}
\label{tab:cost_extended}
\footnotesize
\begin{tabular*}{\columnwidth}{@{\extracolsep{\fill}} l c c c @{}}
\toprule
Comparison & $\Delta$Cost & Window & MIP gap \\
\midrule
L3$\to\Lplus$
  & $+0.11\%$ & $Full year$ & ${<}\,0.5\%$ \\
$\Lplus\to\Lnl$, $Jan.\ 2025$
  & $+0.35\%$ & $744\,h$    & $1.99\%$ \\
$\Lplus\to\Lnl$, $Feb.\ 2025$
  & $+0.50\%$ & $672\,h$    & $0.46\%$ \\
\bottomrule
\end{tabular*}
\end{table}

\paragraph{Linearisation error bounds: $\Lplus$\,$\to$\,$\Lnl$.}
\Cref{tab:cost_extended} reports the combined $\Lplus\to\Lnl$ gap
on two independent winter windows: $+0.35\,\%$ (January, 744\,h)
and $+0.50\,\%$ (February, 672\,h). Three observations bound the
interpretation:
%
\begin{enumerate}[nosep]
  \item The gap combines PWL-vs.-quadratic error
        \emph{and} transport-delay activation; the two cannot be
        separated without an intermediate MIQCP run with delay
        disabled.
  \item The February window has a tight MIP gap (0.46\,\%),
        giving high confidence that the true linearisation effect
        is close to $+0.5\,\%$. The January window has a wider
        gap (1.99\,\%), so its $+0.35\,\%$ is less precisely
        determined.
  \item A spring window (May 2025) was attempted but produced no
        feasible MIQCP incumbent within the time limit because near-zero
        flow in shoulder months causes degenerate bilinear
        constraints. The bound is therefore validated for
        winter conditions (peak-flow regime).
\end{enumerate}
%
The January gap ($+0.35\,\%$) falls below solver tolerance
(MIP gap 1.99\,\%) and is therefore not independently
significant; it does, however, confirm that $\Lnl$ cannot
materially improve on the $\Lplus$ solution even when granted
substantially more optimality slack. The February gap
($+0.50\,\%$, MIP gap 0.46\,\%) is the tighter and
statistically meaningful bound. Together, both windows support
the conclusion that the combined linearisation and
transport-delay effect is ${\leq}\,0.5\,\%$ under winter
operating conditions. The effect is supposedly negligible for planning-accuracy dispatch.

%% ============================================================
%% 5.4 SYNTHETIC PARAMETRIC STUDY
%% ============================================================
\subsection{Generalisability: synthetic parametric study (RQ4)}
\label{sec:results_generalizability}

The synthetic parametric study covers 36~feasible configurations
(\Cref{tab:synthetic}), with the full
L1/L1$_{\mathrm{topo}}$/L2/L3/$\Lplus$ hierarchy solved for each
(180~runs + 4~primary-case = 184 total). Three findings emerge.

\paragraph{Pipe-length scaling.}
The cumulative L1$\to$L3 gap is strictly positive in 36/36
configurations and scales monotonically with pipe length
(\Cref{fig:synth_sweep}c, \Cref{tab:synth_by_length}):
median 3.2\,\% at 1\,\si{\kilo\metre}, 15.8\,\% at
5\,\si{\kilo\metre}, 39.0\,\% at 15\,\si{\kilo\metre}.

\begin{figure*}[!ht]
  \centering
  \includegraphics[width=\textwidth]{paper1_dh_fidelity/F10_synth_sweep_combined.png}
  \caption{Synthetic parametric study across 36~configurations.
  (a)~L1$\to$L3 gap distribution (median 20.0\,\%, range
  3.0--44.8\,\%, 36/36 positive).
  (b)~Additive decomposition sorted by net gap: heat losses
  (red) dominate; topology routing (blue) and pressure drop
  (orange) are negligible.
  (c)~Monotone gap scaling with pipe length.}
  \label{fig:synth_sweep}
\end{figure*}


\begin{table}[ht]
\centering
\caption{Cumulative L1$\to$L3 cost gap by pipe length.}
\label{tab:synth_by_length}
\footnotesize
\begin{tabular*}{\columnwidth}{@{\extracolsep{\fill}} l c c c @{}}
\toprule
Pipe length & $n$ & Median & p10--p90 \\
\midrule
1\,km  &  6 & $+3.2\%$  & $+3.0$--$4.4\%$ \\
5\,km  & 14 & $+15.8\%$ & $+14.2$--$20.1\%$ \\
15\,km & 16 & $+39.0\%$ & $+34.8$--$44.2\%$ \\
\bottomrule
\end{tabular*}
\end{table}

\paragraph{Three-component decomposition.}
Decomposing the gap into topology routing, heat losses, and
pressure drop (\Cref{eq:cost_decomposition}) reveals an
unambiguous driver (\Cref{fig:synth_sweep}b): the heat-loss
component accounts for nearly the entire gap (median
$+20.1\,\%$), while routing ($-0.7$ to $-8\,\%$ of the total
gap\,\footnote{Negative routing contributions indicate that
spatial flow-routing constraints marginally reduce cost by
forcing locally efficient asset utilisation; the heat-loss
component correspondingly exceeds 100\,\% of the net gap.}) and
pressure drop ($+0.02\,\%$) remain at or below the MIP-gap
noise floor. The $+0.02\,\%$ pressure-drop component is
measured within the synthetic decomposition framework
(\Cref{tab:synth_level_mapping}), where L3 includes
Darcy--Weisbach physics; it is consistent with the $+0.11\,\%$
L3$\to\Lplus$ increment for the primary case
(\Cref{tab:cost_extended}), where pressure drop enters one
level later.

\paragraph{Parameter ranking.}
\Cref{tab:synth_marginal} (\Cref{app:synth_detail}) ranks the
four design parameters by their marginal effect on the L1$\to$L3
gap. Pipe length dominates ($+35.8$\,pp across its range);
storage capacity ranks second ($+29.6$\,pp) but is partially
confounded with pipe length due to uneven feasibility filtering
(\Cref{app:synth_detail}). Node count ($+7.1$\,pp) and demand
heterogeneity ($-4.8$\,pp) are an order of magnitude weaker and
non-monotone (\Cref{tab:synth_hi_effect}).

\Cref{fig:synth_surface} confirms this dominance across all
36~configurations: at 1\,km the gap remains below 5\,\%, rising
to 14--20\,\% at 5\,km and exceeding 35\,\% at 15\,km. Storage
capacity has a secondary effect: larger reservoirs allow more
temporal dispatch flexibility that the copperplate model cannot
capture. The nearly vertical gradient confirms that network
topology (via thermal losses) rather than operational scheduling
drives the cost underestimation.

Notably, the L3$\to\Lplus$ gap is identically zero across all
36~synthetic configurations (median and max, for each of the
three pipe-length groups: 1\,km, $n{=}6$; 5\,km, $n{=}14$;
15\,km, $n{=}16$). This
arises because the synthetic L3 implementation includes
Darcy--Weisbach pressure drop (\Cref{tab:synth_level_mapping});
the L3$\to\Lplus$ step therefore isolates only transport-delay
activation, which produces no cost change at hourly resolution
when $\tau_{\max} < 1$\,h. The pressure-drop effect itself
($+0.11\,\%$) is quantified exclusively on the primary case
(\Cref{tab:cost_extended}), where L3 excludes hydraulic physics
by design.

\begin{figure}[!ht]
  \centering
  \includegraphics[width=\columnwidth]{paper1_dh_fidelity/F_synth_gap_surface.png}
  \caption{L1$\to$L3 cost gap as a function of total pipe length
  and thermal storage capacity across all 36~synthetic
  configurations (averaged over node count and demand heterogeneity
  index).}
  \label{fig:synth_surface}
\end{figure}
%% ============================================================
%% 5.5 SENSITIVITY
%% ============================================================
\subsection{Sensitivity analysis}
\label{sec:results_sensitivity}

Gap stability is assessed across 11~price and demand scenarios
(\Cref{tab:gap_stability}). The L1--L3 gap remains in the
12--14\,\% band regardless of fuel (${\pm}25\,\%$), electricity
(${\pm}30\,\%$), or \ce{CO2} price assumptions
(${\pm}50\,\%$), confirming that the topology effect is
\emph{structural} (driven by pipe losses) rather than
\emph{economic} (driven by merit-order shifts). The L3--$\Lplus$
gap stays below 0.3\,\% in all 11~scenarios, reinforcing the
practical irrelevance of extended hydraulic physics.

One notable exception: the warm-year scenario shows a sign
reversal for L2--L3 ($-2.2\,\%$). At reduced thermal demand, pipe losses shrink as a fraction of load; simultaneously, L2's zone-averaged $UL$-product overestimates losses in the warm season, because the spatial demand distribution under partial load deviates most from the calibration baseline. The sign reversal is therefore a loss-aggregation artefact of L2's zone approximation, not a locational asset effect. This is the only scenario in which L2 overestimates cost relative to L3 suggesting that L2 may be
sufficient for low-load or shoulder-season planning even in
heterogeneous networks. 

%% ============================================================
%% 5.6 COMPUTATIONAL PERFORMANCE
%% ============================================================
\subsection{Computational performance}
\label{sec:results_computation}

\Cref{tab:computation} summarises problem size and solve time. All
MILP variants reach the 0.5\,\% gap within 20\,min. $\Lnl$ is
intractable for the full year, confirming $\Lplus$ as the
highest thermo-hydraulic practical formulation.

Storage dispatch is indistinguishable across abstraction levels
on the primary case: cycle counts (8~equivalent winter cycles)
and SOC timing shifts (${<}\,0.2$\,h) are identical across
L1--$\Lplus$. This equivalence is expected for the compact
primary case ($\tau_{\max} \approx 35$\,min); sensitivity to
topology likely emerges in networks with longer transport delays
(\Cref{sec:conclusion}).

\begin{table}[ht]
\centering
\caption{Computational performance (single workstation).}
\label{tab:computation}
\begin{threeparttable}
\scriptsize
\begin{tabular*}{\columnwidth}{@{\extracolsep{\fill}} l r r r r l @{}}
\toprule
Level & Vars & Bin. & Quad.\
  & Solve & Gap \\
\midrule
L1       & 210k   & 53k  & 0   & 268\,s  & ${<}0.5\%$ \\
L2       & 1.28M  & 53k  & 0   & 450\,s  & ${<}0.5\%$ \\
L3       & 2.27M  & 53k  & 0   & 650\,s  & ${<}0.5\%$ \\
$\Lplus$ & 3.02M  & 420k & 0   & 1165\,s & ${<}0.5\%$ \\
$\Lnl$ (744h)
         & --     & --   & 65k & ${>}4$h & 1.99\% \\
\bottomrule
\end{tabular*}
\begin{tablenotes}\scriptsize
  \item $\Lnl$: Jan.\ window, warm-started from $\Lplus$.
\end{tablenotes}
\end{threeparttable}
\end{table}


%% ============================================================
%% 6. DISCUSSION
%% ============================================================
%% ============================================================
%% 6. DISCUSSION
%% ============================================================
\section{Discussion}
\label{sec:discussion}

The results highlight the role of the topology dominance principle in modelling of DH distribution: in radial-tree topologies with fixed supply
temperature and pipe lengths up to 15\,km, the modelling step
that delivers nearly all of the quantified accuracy gain
(spatial resolution with steady-state losses) requires zero
additional solver capability, zero nonlinear constraints, and
increases solve time by merely 70\,\%. Conversely, the modelling
steps that consume the vast majority of development effort and
solver resources (thermo-hydraulic physics with PWL or MIQCP
formulations) deliver less than 1\,\% marginal accuracy. The
practical implication is direct: \emph{resolve the graph before
refining the physics}.

The qualifier matters. The 1\,\% upper bound on extended-physics
effects has been established for radial-tree topologies, fixed
heating curves, and dispatch with given capacities. Ring
topologies with bidirectional flow may amplify hydraulic coupling
through loop-flow degrees of freedom; variable supply temperature
introduces $T \cdot \dot{m}$ bilinear terms that scale with
network throughput; joint dispatch--sizing problems may select
different capacities under L1 vs.\ L3 representations, which in
turn affects pumping fractions. Each of these is identified as a
priority follow-up in \Cref{sec:conclusion}.

\subsection{Effect hierarchy and interpretation}
\label{sec:disc_hierarchy}

The results establish the following magnitude ordering:
%
\begin{align*}
  \underbrace{\text{L1} \to \text{L2}}_{\substack{\text{dominant} \\
    10.4\%}}
  \;>\;
  \underbrace{\text{L2} \to \text{L3}}_{\substack{\text{resolution}
    \\ 2.6\%}}
  \;\gg\;
  \underbrace{\text{L3} \to \Lplus}_{\substack{\Delta p \\
    0.11\%}}
  \;>\;
  \underbrace{\Lplus \to \Lnl}_{\substack{\text{bound} \\
    {\leq}0.5\%}}
\end{align*}
\Cref{fig:effect_hierarchy} visualises the paradox: the first
step (topology + losses) delivers nearly all accuracy gain at
minimal solve-time cost, while subsequent steps invert this
ratio.

\begin{figure}[ht]
  \centering
  \includegraphics[width=\linewidth]{paper1_dh_fidelity/F_effect_hierarchy_combined.png}
  \caption{Effect hierarchy and computational trade-off.
  (a)~Cost gap per modelling step (\% of cost(L3)); topology
  with losses dominates.
  (b)~Additional solve time per step (log scale); extended
  physics and MIQCP consume disproportionate resources for
  marginal accuracy.}
  \label{fig:effect_hierarchy}
\end{figure}

The synthetic decomposition refines this picture: pure topology
routing without losses contributes effectively zero; the entire
dominant step is attributable to thermal losses that become visible
once the network is spatially resolved. An important finding is not that topology matters (this follows from thermodynamics) but that \emph{pressure-drop physics, transport delay, and nonlinear formulations are collectively
unnecessary} for planning-accuracy dispatch in networks up to
15\,km pipe length. This ``negative result'' for thermo-hydraulics has
direct resource-allocation implications: modelling effort should be
invested in spatial resolution (which is cheap computationally) rather
than physics refinement (which is expensive and yields $<\!1\,\%$).

Four implications follow for practitioners:
%
\begin{enumerate}[nosep]
  \item \textbf{Thermal-loss accounting is the first modelling
        investment.} Losses scale with pipe length (median
        $+3\,\%$ at 1\,km to $+39\,\%$ at 15\,km).
  \item \textbf{Pressure-drop physics are marginal.} Even at
        15\,km, the pressure-drop component remains an order of
        magnitude below the topology/loss gap.
  \item \textbf{Transport delay requires MIQCP to characterise.}
        Under MILP linearisation, transport delay is bypassed. If
        $\tau > 1$\,h, a nonlinear validation run is warranted.
  \item \textbf{MILP is adequate for planning.} The combined
        $\Lplus\!\to\!\Lnl$ bound (${\leq}\,0.5\,\%$, February
        window) is itself comparable to solver tolerance; the
        January window confirms no measurable improvement even at
        relaxed optimality (1.99\,\% MIP gap).
\end{enumerate}

\subsection{Integration pathway for multi-energy system frameworks}
\label{sec:disc_tools}

Open-source frameworks for multi-energy system optimisation 
including oemof~\cite{Hilpert2018_oemof},
PyPSA~\cite{Brown2018_pypsa}, FINE~\cite{Welder2018_FINE},
Calliope~\cite{Pfenninger2018_calliope}, and
urbs~\cite{Dorfner2016_urbs}, typically represent thermal
networks at the copperplate level (L1) in their default
DH formulations as used in the MES literature cited above.
Several frameworks offer extensions with per-link loss
representations (e.g.,\ oemof-thermal, PyPSA-Heat, Calliope with loss factors); however, in
cross-framework MES studies, the copperplate representation
remains the de facto default. Our results quantify the cost of
this default: for networks with total pipe length exceeding
\SI{1}{\kilo\metre} (i.e., virtually all operational DH systems),
the L1 default introduces a systematic cost underestimation of
3--45\,\% depending on network extent.

The contribution of this work is therefore to identify a
\emph{pragmatic integration pathway} for frameworks that have not
yet committed to a particular optimisation level:
%
\begin{enumerate}[nosep]
  \item \textbf{L2 as a low-cost minimum.} Adding per-segment
        steady-state losses (\Cref{eq:loss_supply} and \Cref{eq:loss_return}) requires only a linear loss term per
        pipe and a nodal balance reformulation
        (\Cref{eq:balance_nodal}). No solver change, no
        additional binaries, and minimal solve-time increase
        (${+}70\,\%$ in our implementation, \Cref{tab:computation}).
        This step alone closes most of the 3--45\,\% L1-bias gap.
    \item \textbf{L3 covers the main impact: topology} Resolving
        individual junctions (L3) adds modest cost
        (1.44$\times$ L2 solve time) and is warranted when demand
        is spatially heterogeneous (Gini~${>}\,0.3$).
  \item \textbf{Diminishing returns beyond L3.} Pressure-drop
        physics ($\Lplus$) use 420\,k binary variables in total and double
        solve time while contributing ${\leq}\,0.3\,\%$ cost
        accuracy across all tested scenarios (radial trees, fixed
        heating curve; maximum $+0.2\,\%$,
        cf.\ \Cref{tab:gap_stability}). For planning studies, L3
        represents the complexity--accuracy Pareto front.
  \item \textbf{Validation discipline.} Framework users
        integrating network losses should validate against
        measured data using the two-stage protocol described in
        \Cref{sec:val_protocol}; uncalibrated $U$-values can
        introduce systematic errors that exceed the topology gap
        itself.
\end{enumerate}

\Cref{fig:pareto} confirms this hierarchy. In the figure, the solving time is plotted against the accuracy. L3 achieves
${\leq}\,0.5\,\%$ cost deviation from the MIQCP reference at less
than 11\,min solve time, while the full nonlinear reference
($\Lnl$) requires ${>}\,4$\,h for a marginal accuracy gain.

\begin{figure}[ht]
  \centering
  \includegraphics[width=\linewidth]{paper1_dh_fidelity/F_pareto_accuracy_vs_solvetime.png}
  \caption{Accuracy solve-time comparison. The dashed vertical line marks a 1\,h solver run time.}
  \label{fig:pareto}
\end{figure}

These findings do not invalidate L1-based studies for their
stated purpose (national-scale technology screening); they
characterise the magnitude of network-level dispatch error in
default L1 configurations and offer framework developers a
quantitative basis for prioritising thermal-network extensions.

\subsection{Model selection criteria}
\label{sec:disc_criteria}

The quantified gaps translate into a decision tree for
practitioners (\Cref{tab:criteria}). The dominant cost driver is
topology because it is the sole mechanism exposing thermal losses
to the optimiser.

\begin{table*}[!ht]
\raggedright
\caption{Recommended model selection criteria.}
\label{tab:criteria}
\footnotesize
\begin{tabularx}{\textwidth}{@{} l X X l @{}}
\toprule
Level & Conditions & Rationale & Solve \\
\midrule
L1
  & Pipe ${<}\SI{1}{\kilo\metre}$ AND no storage
  & Losses ${<}1\,\%$
  & Seconds \\[4pt]
L2
  & Pipe 1--\SI{10}{\kilo\metre} AND uniform demand
    (Gini ${\leq}0.3$, top-3 ${\leq}50\,\%$)
  & Losses significant; demand uniform
  & Minutes \\[4pt]
L3
  & Heterogeneous demand
    (Gini ${>}0.3$ or top-3 ${>}50\,\%$) OR
    storage ${>}6$\,h
  & Spatial dispatch differences
  & Min- to hours \\[4pt]
$\Lplus$
  & Pipe ${>}\SI{5}{\kilo\metre}$ AND pumping ${>}1\,\%$
    of thermal losses
  & Extended hydraulic physics needed
  & Hours \\[4pt]
$\Lnl$
  & $\tau_{\max} > \SI{1}{\hour}$ (transport delay relevant);
    OR validation of $\Lplus$;
    OR research
  & Only level with active transport delay
  & Hours to days \\
\bottomrule
\end{tabularx}
\end{table*}
\subsection{\ce{CO2} policy implications}
\label{sec:disc_co2}

Simplified models underestimate \ce{CO2} through compounding
channels: missed loss-compensation generation, missed pumping
electricity, and temporal correlation between additional demand
and grid emission intensity. In the primary case, L1-based
accounting underestimates emissions by $8.7\,\%$ (772
vs.\ 846\,t\,\ce{CO2}/yr); the synthetic parametric study shows this scaling
from ${<}5\,\%$ at 1\,km to ${>}30\,\%$ at 15\,km.

These results suggest that \ce{CO2} reporting frameworks applied
to district heating may need to specify a minimum
network-resolution level. Based on the present evidence, L2
resolution appears warranted whenever the pipe length exceeds
1--2\,km, though this recommendation rests on a single industrial
case and a synthetic parametric study covering radial-tree topologies with
fixed heating curves. Broader empirical validation across diverse
network types is needed before this finding can support binding
regulatory specifications.

\subsection{Limitations}
\label{sec:discussion_limitations}

\paragraph{Validation and generalisation.}
Direct validation covers the pre-upgrade network (Stage~1);
post-upgrade HP/TES dispatch is verified through the necessary
conditions only. The 13\,\% headline gap is from one industrial
network; the synthetic parametric study provides qualitative external validity,
but absolute figures should not be treated as universal constants.

\paragraph{Deterministic price and demand assumptions.}
Dispatch here is evaluated under deterministic price and demand
scenarios. \citet{Koek2026} show, for stochastic multistage
investment planning of DH generation capacity, that ignoring price
and demand-trajectory uncertainty (the Value of the Stochastic
Solution) biases technology selection by 27--39\,\%; while their
setting is investment planning rather than dispatch, and network
topology is not a variable in their framework, the result suggests
that the topology bias quantified here and the uncertainty-driven
bias they quantify are likely additive rather than substitutable,
and jointly addressing both remains an open extension.

\paragraph{Scope restrictions.}
The framework is dispatch-only (fixed capacities); joint
dispatch-and-sizing may interact with topology choice. Supply
temperature follows a heating curve (not optimised) to preserve
MILP linearity; variable-temperature operation could amplify the
L3$\to\Lplus$ effect. \citet{Wirtz2021_tempcontrol} show, for a
5GDHC network, that co-optimizing network temperature within an
MILP via temperature-level discretisation is tractable and
cost-reducing relative to free-floating temperature control;
applying this discretisation approach within the present
topology-isolated framework is a natural extension.

\paragraph{Demand concentration metric.}
The normalised variance index HI (defined in
\Cref{app:hi_definition}) is used as a parametric study variable but has a
narrow dynamic range on realistic tree networks (HI\,=\,0.05 at
Gini\,=\,0.41 for the primary case). The selection criteria
(\Cref{tab:criteria}), therefore, use Gini and top-$k$ share. An
empirical mapping from these metrics to the L2$\to$L3 gap is a
natural follow-up requiring multiple real networks.

\paragraph{$\Lplus\!\to\!\Lnl$ interpretation.}
The $+0.35$ to $+0.50\,\%$ gap across two independent winter
windows combines the linearisation error and transport-delay
activation; the two contributions cannot be separated without an
intermediate MIQCP run with delay disabled. The bound is,
therefore, an upper limit on the combined effect, not a precise
decomposition.

\paragraph{$\Lnl$ tractability.}
The MIQCP reference is solved on two winter windows (October to
February); a spring window (May) proved infeasible due to
near-zero-flow degeneracy in bilinear constraints. The bound is
therefore validated for peak-flow winter conditions. Whether the
linearisation error remains sub-percent under shoulder-season
operating points with frequent on/off cycling is an open question
best addressed through warm-start heuristics or reformulation.

\paragraph{Demand-side flexibility.}
All demands are exogenous time series; demand response is
excluded. If flexible loads are modelled, the topology effect may
increase because spatially distributed flexibility has
location-dependent value (a flexible load near j$_{12}$ can
absorb local HP surplus without trunk losses). This interaction
between demand flexibility and topology is relevant for real-time
dispatch.

\paragraph{Network topology class.}
All cases (primary and synthetic) are radial trees. Ring topologies
with bidirectional flow may amplify hydraulic coupling effects and
shift the L3$\to\Lplus$ threshold; this remains an open question.
%% ============================================================
%% 7. CONCLUSION
%% ============================================================
\section{Conclusion}
\label{sec:conclusion}

This paper presents a comparison that systematically
isolates topology abstraction, physics fidelity, and linearisation
error as independent experimental variables in district heating
dispatch optimisation. This study varied three dimensions, one at a time, within a
single experimental framework.

\paragraph{Principal findings}
\begin{enumerate}[nosep]
  \item \textbf{Topology effect (RQ1):} The cumulative L1$\to$L3
        gap reaches 13.0\,\% in cost and 74\,t\,\ce{CO2}/year.
        Across 36~synthetic networks, it scales monotonically
        with pipe length (3.2\,\% at 1\,km, 15.8\,\% at 5\,km,
        39.0\,\% at 15\,km; 36/36 positive).
  \item \textbf{Physical driver:} Thermal losses account for
        nearly the entire effect (median $+20.1\,\%$); pure
        routing ($-0.2\,\%$) and pressure drop ($+0.02\,\%$)
        are negligible. The key insight is that topology matters
        not because of routing constraints but because it is the
        only mechanism making pipe losses visible to the optimiser.
  \item \textbf{Physics fidelity (RQ2):}
        Extended hydraulic and thermal physics do not materially
        change dispatch when topology is already resolved:
        L3$\to\Lplus$ adds $+0.11\,\%$, attributable entirely
        to pumping power.
   \item \textbf{Linearisation and transport delay (RQ3):}
        The combined $\Lplus\to\Lnl$ gap is bounded at
        ${\leq}\,0.5\,\%$ based on the February window
        ($+0.50\,\%$, MIP gap 0.46\,\%); the January window
        ($+0.35\,\%$) is consistent but falls below solver
        tolerance (MIP gap 1.99\,\%) and is therefore not
        independently significant. Both windows confirm that
        $\Lnl$ cannot materially outperform $\Lplus$, supporting
        MILP as adequate for planning-accuracy dispatch. A precise
        separation of linearisation error from transport-delay
        activation remains open.
  \item \textbf{Generalisability (RQ4):} Pipe length is the
        dominant predictor; other parameters show weak
        within-band correlations.
\end{enumerate}

\paragraph{Recommendations}
\begin{itemize}[nosep]
  \item L2 as minimum for pipe ${>}\SI{1}{\kilo\metre}$ or
        storage present.
  \item L3 for spatially heterogeneous networks (Gini ${>}0.3$).
  \item $\Lplus$ when pipe ${>}\SI{5}{\kilo\metre}$ or
        $\tau > \SI{1}{\hour}$.
  \item $\Lnl$ as a one-time validation benchmark.
\end{itemize}

\paragraph{Future work}
The research identified several continues research topics: storage-dispatch sensitivity on networks with longer transport delays; ring and meshed topologies; variable transport delay; supply-temperature optimisation ; joint dispatch-sizing comparison; full MINLP with native exponentials; stochastic variants; sub-hourly resolution; and 5th-generation bidirectional networks. Independent separation of the linearisation error from transport-delay activation via intermediate MIQCP runs is a
priority follow-up.

\paragraph{Closing remark}
A single-node model underestimates annual operating costs by
13\,\% and \ce{CO2} emissions by 8.7\,\% on the primary
case, a gap that scales from 3\,\% to 45\,\% across the
36~synthetic configurations, driven almost entirely by thermal losses made
visible through spatial resolution. Within the tested scope
(radial trees, fixed heating curve, dispatch-only), extended
hydraulic and thermal physics add less than 1\,\% on top, and
MILP linearisation introduces a further sub-percent margin on the
validation windows. The practical implication for this regime is
direct: invest modelling effort in spatial resolution rather than
physics refinement. Whether the same hierarchy holds for ring
topologies, variable supply temperature, or joint
dispatch--sizing problems is a question for the follow-up
studies identified above.

%% ============================================================
%% ACKNOWLEDGMENTS
%% ============================================================
\section*{Acknowledgments}
This work was funded by the German Federal Ministry for Economic
Affairs and Energy (BMWE, reference IIC2) under grant
number~\textbf{03EN6057B} (Verbundvorhaben eProNe; project period 2026--2028).
Project management was carried out by Forschungszentrum J\"ulich
GmbH (PT-J.ESN6). The authors thank the e-con AG in
Memmingen for providing operational data and domain expertise.

%% ============================================================
%% DATA AVAILABILITY
%% ============================================================
\section*{Data availability}

The optimisation  model source code (Python/Pyomo), all YAML
configuration files, and post-processing scripts are available on
Zenodo at \citet{ruess_2026_calion} under
"Creative Commons Attribution 4.0 International". Input time series are subject to NDA; anonymised summary
statistics are included in the repository.

%% ============================================================
%% AI Use
%% ============================================================
\section*{Declaration of generative AI and AI-assisted technologies in the manuscript preparation process}

During the preparation of this work, the author(s) used
\textbf{Claude Opus 4.7} (Anthropic) for language improvement and
manuscript structuring, and \textbf{Claude Sonnet 4.6} (Anthropic)
for code generation and debugging.  All content was reviewed and
edited by the authors, who take full responsibility for the
published article.

%% ============================================================
%% APPENDICES
%% ============================================================

\appendix
\newpage

\section{Nomenclature}
\label{app:nomenclature}

{\footnotesize

%%--- Sets and indices ---
\noindent\textbf{Sets and indices}\par\nopagebreak\smallskip
\noindent
\begin{tabular*}{\columnwidth}{@{\extracolsep{\fill}} >{\raggedright}p{2.0cm} l @{}}
\toprule
Symbol & Description \\
\midrule
$t \in \mathcal{T}$                & Time steps ($|\mathcal{T}|=8760$) \\
$n \in \mathcal{N}$                & Network nodes \\
$p \in \mathcal{P}$                & Pipe segments \\
$g \in \mathcal{G}$                & Thermal generators \\
$h \in \mathcal{H}$                & Heat pumps \\
$s \in \mathcal{S}$                & Thermal storage units \\
$r \in \mathcal{R}$                & Electrode boilers (power-to-heat) \\
$k \in \{1,\ldots,K\}$            & PWL segment index \\
$\mathcal{G}^{\mathrm{CHP}}$      & Subset of CHP generators \\
$\mathcal{P}_n^{\mathrm{in/out}}$  & Pipes entering/leaving node~$n$ \\
$\mathcal{G}_n,\mathcal{S}_n$     & Generators/storage at node~$n$ \\
$\mathcal{T}_{\mathrm{val}}$      & Validation time steps \\
\bottomrule
\end{tabular*}

\bigskip

%%--- Key parameters ---
\noindent\textbf{Key parameters}\par\nopagebreak\smallskip
\noindent
\begin{tabular*}{\columnwidth}{@{\extracolsep{\fill}} >{\raggedright}p{2.0cm} >{\raggedright}p{2.0cm} l @{}}
\toprule
Symbol & Unit & Description \\
\midrule
$U_p$                        & W\,m$^{-1}$\,K$^{-1}$ & Pipe heat transfer coeff.\ \\
$L_p$                        & m              & Pipe length \\
$D_p$                        & m              & Pipe inner diameter \\
$A_p$                        & m$^2$          & Pipe cross-sectional area \\
$R_p$                        & Pa\,MW$^{-2}$  & Pipe hydraulic resistance \\
$f_D$                        & --             & Darcy friction factor \\
$\rho$                       & kg\,m$^{-3}$   & Fluid density \\
$c_p$                        & J\,(kg\,K)$^{-1}$ & Specific heat capacity \\
$\tau_p$                     & h              & Transport delay \\
$k_p$                        & --             & Delay in time steps \\
$\phi_{p,t}$                 & --             & Temp.\ decay factor \\
$\phi_p^{\mathrm{nom}}$      & --             & Nominal decay factor \\
$T_{\mathrm{sup}}^{\mathrm{nom}}(t)$ & $^\circ$C & Nominal supply temp.\ \\
$T_{\mathrm{ret}}^{\mathrm{nom}}$    & $^\circ$C & Nominal return temp.\ \\
$T^{\mathrm{nom}}(t)$        & $^\circ$C      & Nominal operating temp.\ \\
$T_{\mathrm{ground}}(t)$     & $^\circ$C      & Ground temperature \\
$\bar{Q}_g,\bar{Q}_h,\bar{Q}_r$ & MW         & Installed thermal capacities \\
$\bar{Q}_p$                  & MW             & Rated pipe capacity \\
$\phi_h^{\mathrm{min}}$      & --             & HP minimum load fraction \\
$\eta_g$                     & --             & Generator thermal efficiency \\
$\eta_s^{\mathrm{ch/dis}}$   & --             & Storage charge/discharge eff.\ \\
$\eta_{\mathrm{pump}}$       & --             & Pump efficiency \\
$\eta_{\mathrm{Lorenz}}$     & --             & Lorenz cycle efficiency \\
$\alpha_s$                   & h$^{-1}$       & Storage self-discharge rate \\
$\mathrm{COP}_{h,t}^{\mathrm{wrg}}$ & --     & Waste-heat COP (time-varying) \\
$\mathrm{COP}_h^{\mathrm{def}}$     & --     & Default-source COP (constant) \\
$m_k^{(p)}$                  & --             & PWL segment slope \\
$\Delta T_{\mathrm{pp}}$     & K              & Heat pump pinch-point temp.\ \\
$\lambda_t^{\mathrm{buy/sell}}$ & \euro\,MWh$^{-1}$ & Electricity buy/sell price \\
$\lambda_g^{\mathrm{fuel}}$  & \euro\,MWh$^{-1}$ & Specific fuel cost \\
$\lambda^{\mathrm{dump}}$    & \euro\,MWh$^{-1}$ & Heat curtailment penalty \\
$\lambda^{\ce{CO2}}$         & \euro\,t$^{-1}$ & Carbon price \\
$\lambda^{\mathrm{dem}}$     & \euro\,MW$^{-1}$ & Demand charge rate \\
$e_t^{\mathrm{grid}}$        & kg\,MWh$^{-1}$ & Grid emission factor \\
$e_g^{\mathrm{fuel}}$        & kg\,MWh$^{-1}$ & Fuel emission intensity \\
$f^{\mathrm{yr}}$            & --             & Annual scaling factor \\
$M^{\mathrm{grid}}$          & MW             & Grid capacity (Big-M) \\
$\Delta t$                   & h              & Time step length ($= 1$\,h) \\
HI                           & --             & Demand heterog.\ index \\
$K$                          & --             & Number of PWL segments \\
\bottomrule
\end{tabular*}

\bigskip

%%--- Cost components ---
\noindent\textbf{Cost components}\par\nopagebreak\smallskip
\noindent
\begin{tabular*}{\columnwidth}{@{\extracolsep{\fill}} >{\raggedright}p{2.0cm} l @{}}
\toprule
Symbol & Description \\
\midrule
$C^{\mathrm{en}}$    & Electricity procurement net cost \\
$C^{\mathrm{fuel}}$  & Fuel expenditure \\
$C^{\mathrm{dump}}$  & Heat curtailment penalty \\
$C^{\ce{CO2}}$       & Carbon cost \\
$C^{\mathrm{dem}}$   & Demand charge (peak import) \\
$C^{\mathrm{pump}}$  & Network pumping cost \\
$Z$                  & Total annual operational cost \\
\bottomrule
\end{tabular*}

\bigskip

%%--- Key decision variables ---
\noindent\textbf{Key decision variables}\par\nopagebreak\smallskip
\noindent
\begin{tabular*}{\columnwidth}{@{\extracolsep{\fill}} >{\raggedright}p{2.0cm} >{\raggedright}p{2.0cm} l @{}}
\toprule
Symbol & Unit & Description \\
\midrule
$Q_{g,t}^{\mathrm{th}}$       & MW       & Generator thermal output \\
$Q_{h,t}$                     & MW       & Heat pump thermal output \\
$Q_{h,t}^{\mathrm{wrg/def}}$  & MW       & HP waste-heat/default output \\
$Q_{r,t}$                     & MW       & Electrode boiler output \\
$Q_{p,t}$                     & MW       & Pipe heat flow \\
$Q_{p,t}^{\mathrm{loss,pipe}}$& MW       & Pipe thermal loss \\
$Q_{s,t}^{\mathrm{ch/dis}}$   & MW       & Storage charge/discharge rate \\
$Q_{n,t}^{\mathrm{dem}}$      & MW       & Nodal heat demand \\
$Q_{n,t}^{\mathrm{dump}}$     & MW       & Curtailed heat at node~$n$ \\
$P_t^{\mathrm{buy/sell}}$     & MW       & Grid import/export \\
$P_{g,t}^{\mathrm{el}}$       & MW       & CHP electrical output \\
$P_{h,t}^{\mathrm{el}}$       & MW       & Heat pump electrical input \\
$P_{r,t}^{\mathrm{el}}$       & MW       & Electrode boiler electrical input \\
$P_t^{\mathrm{pump}}$         & MW       & Network pumping power \\
$P^{\mathrm{peak}}$           & MW       & Annual peak grid import \\
$\Delta p_{p,t}$              & Pa       & Pipe pressure drop \\
$\dot{m}_{p,t}$               & kg\,s$^{-1}$ & Pipe mass flow rate \\
$T_{n,t}^{\mathrm{sup}}$      & $^\circ$C & Node supply temperature \\
$T_{n_1,t}^{\mathrm{sup}}$    & $^\circ$C & Upstream node supply temp.\ \\
$T_{p,t}^{\mathrm{out}}$      & $^\circ$C & Pipe outlet temperature \\
$E_{s,t}$                     & MWh      & Storage state of charge \\
$F_{g,t}$                     & MW       & Fuel input (HHV basis) \\
$E_t^{\ce{CO2},\mathrm{grid}}$ & t       & Grid \ce{CO2} emissions \\
$E_{g,t}^{\ce{CO2},\mathrm{fuel}}$ & t   & Fuel \ce{CO2} emissions \\
$W_{p,t}$                     & MW$^2$   & Bilinear auxiliary ($= Q_{p,t}^2$) \\
$u_{h,t}$                     & $\{0,1\}$ & Heat pump on/off \\
$z_t$                         & $\{0,1\}$ & Grid import mode \\
$w_{p,t,k}$                   & $[0,1]$  & SOS2 weights \\
\bottomrule
\end{tabular*}

}% end footnotesize

%% ============================================================
\section{Model formulation details}
\label{app:model_details}

\subsection{COP pre-computation}
\label{app:cop}

The heat pump COP is pre-computed from an analytical Lorenz-cycle model:
\begin{equation}
  \mathrm{COP}_{h,t}^{\mathrm{WHR}} = \eta_{\mathrm{Lorenz}}
  \cdot \frac{T_{\mathrm{sink,out}}}{T_{\mathrm{sink,out}} -
  T_{\mathrm{src},h}(t) + \Delta T_{\mathrm{pp}}}
\end{equation}
with $\eta_{\mathrm{Lorenz}} = 0.6$ and
$\Delta T_{\mathrm{pp}} = \SI{5}{\kelvin}$.  Combined error:
$\varepsilon_{\mathrm{COP}} \approx 3.6\,\%$.

\subsection{Linear component models}
\label{app:component_models}

{%
\captionof{table}{Linear component models (all levels).}
\label{tab:simple_components}
\centering
\footnotesize
\begin{tabular*}{\columnwidth}{@{\extracolsep{\fill}} l l l @{}}
\toprule
Component & Equation & Bound \\
\midrule
Gas boiler    & $Q_{g,t}^{\mathrm{th}} = \eta_g F_{g,t}$ & $\leq \bar{Q}_g$ \\
Biomass       & $Q_{g,t}^{\mathrm{th}} = \eta_g F_{g,t}$ & $\leq \bar{Q}_g$ \\
Electrode     & $Q_{r,t} = \eta_r P_{r,t}^{\mathrm{el}}$ & $\leq \bar{Q}_r$ \\
CHP           & $Q^{\mathrm{th}} = \eta^{\mathrm{th}} F$  & $\sigma_g = \mathrm{const}$ \\
\bottomrule
\end{tabular*}
}% end centering group

\bigskip

\subsection{Extended physics: $\Lplus$ vs.\ $\Lnl$}
\label{app:extended_formulations}

{%
\captionof{table}{Mathematical treatment of extended physics.}
\label{tab:lplus_vs_lnl}
\centering
\footnotesize
\begin{threeparttable}
\begin{tabular*}{\columnwidth}{@{\extracolsep{\fill}} l l l @{}}
\toprule
Phenomenon & $\Lplus$ & $\Lnl$ \\
\midrule
Pressure drop   & PWL ($K{=}3$)     & Quad.\ $RQ^2$ \\
Pumping power   & PWL of cubic      & Bilin.\ $Q{\cdot}W$ \\
Temp.\ decay    & PWL exp.\tnote{*} & PWL exp.\tnote{*} \\
Temp.\ mixing   & Taylor lin.       & Bilin.\ $T{\cdot}\phi$ \\
Transport delay & Bypassed          & Active ($k_p$) \\
Optimality      & Global (LP)       & Global (branch.) \\
\bottomrule
\end{tabular*}
\begin{tablenotes}\scriptsize
  \item[*] Shared; excluded from linearisation error.
\end{tablenotes}
\end{threeparttable}
}% end centering group

\bigskip

\subsection{Linearisation error bounds}
\label{app:pwl}

For $\Delta p = R Q^2$ with $K$ equal PWL segments:
\begin{equation}
  \varepsilon_{\Delta p}^{\max} = \frac{R \bar{Q}^2}{4K^2}
\end{equation}
With $K{=}3$: 2.8\,\% of peak pressure drop. The cost gap
$\Delta Z$ is well below $\varepsilon^{\max}$ because the
optimiser compensates by adjusting dispatch rather than by
violating feasibility; the PWL error thus manifests as a minor
scheduling shift, not as a proportional cost deviation.

\medskip
{%
\captionof{table}{PWL approximation error and problem size
vs.\ number of segments $K$ (primary case).}
\label{tab:pwl_sensitivity}
\centering
\footnotesize
\begin{threeparttable}
\begin{tabular*}{\columnwidth}{@{\extracolsep{\fill}} c c r l @{}}
\toprule
$K$ & $\varepsilon^{\max}$ & Vars\tnote{a} & $\Delta Z$ \\
\midrule
3 & 2.8\,\% & 753k  & $+0.35\,\%$ \\
5 & 1.0\,\% & 998k\tnote{b}  & -- \\
8 & 0.4\,\% & 1366k\tnote{b} & -- \\
\bottomrule
\end{tabular*}
\begin{tablenotes}\scriptsize
  \item[a] Total additional variables relative to L3
    (binary + continuous).
  \item[b] Extrapolated from the per-segment increment
    (${\approx}\,122$k variables per additional segment);
    not simulated.
\end{tablenotes}
\end{threeparttable}
}% end centering group


\subsection{Taylor linearisation}
\label{app:taylor_derivation}

The bilinear product $T_{n_1} \cdot \phi_p$ is expanded around
$(T_{n_1}^0, \phi_p^0)$:
\begin{equation}
  T_{n_1} \cdot \phi_p \approx T_{n_1}^0 \phi_p^0
  + \phi_p^0 (T_{n_1} - T_{n_1}^0)
  + T_{n_1}^0 (\phi_p - \phi_p^0)
\end{equation}
Neglected cross-term:
$|\varepsilon| \leq |\Delta T_{\max}| \cdot |\Delta\phi_{\max}| < 2.5\,$K.

\subsection{Electricity selling price}
\label{app:selling_price}

\begin{equation}
  \lambda_t^{\mathrm{sell}} = \max\!\Bigl(
    \max\bigl(\lambda_t^{\mathrm{spot}}
      - \lambda^{\mathrm{spread}},\;
      \lambda^{\mathrm{floor}}\bigr)
    (1 - h^{\mathrm{sell}})
    - \lambda^{\mathrm{sell,fee}},\; 0\Bigr)
\end{equation}

\subsection{Per-node MAE breakdown}
\label{app:per_node_mae}

\Cref{tab:per_node_mae} reports supply-temperature MAE and bias
for all 14~network nodes over the 744\,h winter validation window,
sorted by path length from the source.  Three node groups are
distinguished:
%
\begin{itemize}[nosep]
  \item \textbf{Validation nodes} (val., 6~nodes): consumer sensors
        with mixing-valve offset ${<}\SI{2}{\celsius}$,
        independent of the calibration procedure.  Mean MAE
        \SI{1.32}{\celsius}; these are the nodes that enter the
        gate evaluation in \Cref{tab:validation_stage1}.
  \item \textbf{Calibration nodes} (cal., 4~nodes): used in the
        branch-sequential $U_p$-tuning of
        \citet{Maldonado2024}; their residuals are in-sample by
        construction.  Mean MAE \SI{2.32}{\celsius}.  Node j$_7$ at
        \SI{3.07}{\celsius} would fail the worst-node gate of
        \SI{2.5}{\celsius} if held out for validation.  This
        in-sample deterioration is expected and motivates the
        strict calibration/validation split.
  \item \textbf{Excluded nodes} (excl., 4~nodes): excluded for
        measurement reasons unrelated to model quality.  Nodes
        j$_3$ (\SI{0.58}{\celsius}) and j$_4$
        (\SI{0.37}{\celsius}) are the two best-performing nodes
        in the entire network; j$_4$ has four consumers with a
        \SI{20.9}{\celsius} intra-node spread, so no single
        junction-reference sensor exists.  j$_3$ is a
        boundary-condition reference.  j$_2$ and j$_5$ exhibit
        mixing-valve absorption above the \SI{2}{\celsius}
        threshold and are therefore not interpretable as junction
        temperatures.
\end{itemize}
%
The all-14-node mean MAE is \SI{1.68}{\celsius}, which sits
between the validation-node mean (\SI{1.32}{\celsius}) and the
calibration-node mean (\SI{2.32}{\celsius}) and confirms that the
model performance is not artificially inflated by selective node
choice: the validation nodes are not systematically the
best-performing ones.

\medskip
{%
\captionof{table}{Per-node supply-temperature MAE and bias
(744\,h winter window, $\Lnl$).}
\label{tab:per_node_mae}
\centering
\footnotesize
\begin{threeparttable}
\begin{tabular*}{\columnwidth}{@{\extracolsep{\fill}} l c c c c c @{}}
\toprule
Node & Path [m] & Role & MAE [$^\circ$C] & Bias [$^\circ$C] & In gate? \\
\midrule
j$_2$    &  350 & excl. & 2.99 & $+2.98$ & No  \\
j$_3$    &  650 & excl. & 0.58 & $+0.07$ & No  \\
j$_4$    &  910 & excl. & 0.37 & $-0.18$ & No  \\
j$_9$    & 1100 & val.  & 0.87 & $+0.41$ & Yes \\
j$_5$    & 1150 & excl. & 2.35 & $+1.51$ & No  \\
j$_6$    & 1300 & cal.  & 1.62 & $-1.62$ & No  \\
j$_{10}$ & 1300 & val.  & 2.27 & $+1.78$ & Yes \\
j$_7$    & 1370 & cal.  & 3.07\tnote{a} & $+2.53$ & No  \\
j$_8$    & 1450 & cal.  & 2.29 & $+1.00$ & No  \\
j$_{11}$ & 1530 & val.  & 1.41 & $+0.71$ & Yes \\
j$_{12}$ & 1780 & val.  & 0.69 & $-0.29$ & Yes \\
j$_{13}$ & 2000 & val.  & 1.09 & $-1.02$ & Yes \\
j$_{15}$ & 2125 & val.  & 1.58 & $-0.24$ & Yes \\
j$_{14}$ & 2180 & cal.  & 2.31 & $+1.02$ & No  \\
\midrule
\multicolumn{3}{@{}l}{Group mean, val.\ (6~nodes)}  & 1.32 & & \\
\multicolumn{3}{@{}l}{Group mean, cal.\ (4~nodes)}  & 2.32 & & \\
\multicolumn{3}{@{}l}{Group mean, excl.\ (4~nodes)} & 1.57 & & \\
\multicolumn{3}{@{}l}{All 14~nodes}                  & 1.68 & & \\
\bottomrule
\end{tabular*}
\begin{tablenotes}\scriptsize
  \item[a] In-sample residual (calibration node); would exceed
    the worst-node validation gate of \SI{2.5}{\celsius}
    if held out for validation.
\end{tablenotes}
\end{threeparttable}
}% end centering group

\bigskip

\section{Demand heterogeneity index}
\label{app:hi_definition}

The normalised variance index HI is defined as:
\begin{equation}
  d_n = \frac{\sum_t Q_{n,t}^{\mathrm{dem}}}{\sum_n \sum_t
    Q_{n,t}^{\mathrm{dem}}},
  \quad
  \mathrm{HI} = \frac{N}{N{-}1} \sum_{n=1}^{N}
    \left(d_n - \frac{1}{N}\right)^{\!2}
  \label{eq:hi}
\end{equation}
HI maps monotonically to Gini for tree topologies but compresses
the realistic range (HI\,$= 0.05$ for the primary case despite
Gini\,$= 0.41$).  The synthetic-study values
(HI\,$= 0.1, 0.4, 0.8$) correspond approximately to Gini ranges
of 0.15--0.25, 0.40--0.55, and 0.70--0.85.

\section{BCM calibration details}
\label{app:bcm_calibration}

The BCM forward reconstruction distributes heat loss uniformly
across the full \SI{2125}{\metre} trunk corridor via a global
multiplier (${\times}1.330$ on nominal EN~253 $U$-values for all
8~trunk pipes).  The multiplier was determined by golden-section
search, minimising MAE on the Oct--Feb winter subset.  This
approach differs from the branch-sequential optimisation -model
calibration (which concentrates losses in the terminal pipe at
$4.7{\times}$~nominal); both methods are valid for their
respective purposes (BCM: forward physics; L3/$\Lnl$: dispatch
optimisation ).  RMSE: \SI{2.15}{\celsius}; bias:
${+}\SI{0.72}{\celsius}$; 60\,\% of hourly steps within
${\pm}\SI{1.5}{\celsius}$.

%% ============================================================
\section{Synthetic parametric study: detailed parameter effects}
\label{app:synth_detail}

The following tables provide the quantitative evidence underlying
the parameter-ranking statements in
\Cref{sec:results_generalizability}.

\paragraph{Implementation note.}
In the synthetic parametric study, the L3 configuration includes
steady-state pressure-drop calculation (Darcy--Weisbach) as part
of the pipe model, unlike the primary-case L3, which excludes
pressure drop (\Cref{tab:model_scope}). This reassignment is
required by the additive decomposition
(\Cref{eq:cost_decomposition,tab:synth_level_mapping}).
Consequently, the L3$\to\Lplus$ increment in the synthetic study
isolates \emph{only} transport-delay activation, which produces
zero cost difference at hourly resolution with
$\tau_{\max} < 1$\,h. The primary-case L3$\to\Lplus$ gap
($+0.11\,\%$, \Cref{tab:cost_extended}) remains the valid
reference for the isolated pressure-drop effect.

%%--- Tab A: HI effect within pipe-length bands ---
\medskip
{%
\captionof{table}{L1$\to$L3 cost gap [\%] by pipe length and
demand heterogeneity index (median over node count and storage
capacity).}
\label{tab:synth_hi_effect}
\centering
\begin{threeparttable}
\footnotesize
\begin{tabular*}{\columnwidth}{@{\extracolsep{\fill}} l c c c c @{}}
\toprule
Pipe [km] & HI\,=\,0.1 & HI\,=\,0.4\tnote{a} & HI\,=\,0.8
  & $\Delta$(High--Low) \\
\midrule
1   & $+3.23$ & -- & $+3.19$ & $-0.04$\,pp \\
5   & $+15.88$ & -- & $+14.55$ & $-1.33$\,pp \\
15  & $+38.46$ & -- & $+38.63$ & $+0.17$\,pp \\
\bottomrule
\end{tabular*}
\begin{tablenotes}\scriptsize
  \item[a] Median not reported where $n < 3$ feasible
    configurations per cell.
\end{tablenotes}
\end{threeparttable}
}% end group

\bigskip

%%--- Tab B: Decomposition extremes ---
{%
\captionof{table}{Additive gap decomposition for the three
largest-gap and three smallest-gap configurations.  Heat losses
account for 100--108\,\% of the net gap; routing is slightly
negative in all cases.}
\label{tab:synth_decomp_extremes}
\centering
\footnotesize
\setlength{\tabcolsep}{2pt}
\begin{tabular*}{\columnwidth}{@{\extracolsep{\fill}}
  c c c c c c c c @{}}
\toprule
Rank & Pipe & HI & Nodes & Stor.
  & Routing & Loss & $\Delta P$ \\
  & [km] & & & [h] & [\%] & [\%] & [\%] \\
\midrule
\multicolumn{8}{@{}l}{\textit{Top-3 (largest total gap)}} \\
1 & 15 & 0.1 & 30 & 12 & $-$3.5 & $+$47.2 & $+$1.1 \\
2 & 15 & 0.8 & 30 & 12 & $-$4.1 & $+$46.8 & $+$0.9 \\
3 & 15 & 0.1 & 15 & 12 & $-$2.8 & $+$44.1 & $+$1.5 \\
\midrule
\multicolumn{8}{@{}l}{\textit{Bottom-3 (smallest total gap)}} \\
34 & 1 & 0.8 & 5  & 6  & $-$0.2 & $+$3.3 & $+$0.0 \\
35 & 1 & 0.1 & 5  & 2  & $-$0.1 & $+$3.1 & $+$0.0 \\
36 & 1 & 0.1 & 5  & 6  & $-$0.1 & $+$3.0 & $+$0.0 \\
\bottomrule
\end{tabular*}
}% end group

\bigskip


\bigskip

%%--- Tab D: Parameter ranking ---
{%
\captionof{table}{Marginal effect of each design parameter on
the L1$\to$L3 cost gap (median gap at each parameter level,
averaged over remaining parameters; sorted by effect size).}
\label{tab:synth_marginal}
\centering
\begin{threeparttable}
\footnotesize
\setlength{\tabcolsep}{3pt}
\begin{tabular*}{\columnwidth}{@{\extracolsep{\fill}}
  c
  >{\raggedright\arraybackslash}p{0.20\columnwidth}
  >{\centering\arraybackslash}p{0.11\columnwidth}
  >{\centering\arraybackslash}p{0.11\columnwidth}
  >{\centering\arraybackslash}p{0.11\columnwidth}
  >{\centering\arraybackslash}p{0.14\columnwidth} @{}}
\toprule
Rank & Parameter & Low & Med & High
  & $\Delta$\,(H--L) \\
\midrule
1 & Pipe length [km]
  & $+3.2\,\%$ & $+15.8\,\%$ & $+39.0\,\%$
  & $+35.8$\,pp \\
2 & Storage cap.\ [h]\tnote{a}
  & $+9.4\,\%$ & $+17.2\,\%$ & $+39.0\,\%$
  & $+29.6$\,pp \\
3 & Node count
  & $+12.9\,\%$ & $+16.8\,\%$ & $+20.0\,\%$
  & $+7.1$\,pp \\
4 & Demand HI
  & $+20.8\,\%$ & -- & $+16.0\,\%$
  & $-4.8$\,pp \\
\bottomrule
\end{tabular*}
\begin{tablenotes}\scriptsize
  \item[a] Storage-capacity effect partially confounded with
    pipe length: 2\,h configurations are unevenly distributed
    across pipe-length bands (4/6 occur at $L \geq 5$\,km).
    The isolated storage effect within fixed pipe-length bands
    is smaller (${\approx}5$--$8$\,pp).
\end{tablenotes}
\end{threeparttable}
}% end group
%% ============================================================
\section{Case study supplementary data}
\label{app:case_data}

\subsection{Heating curve}
\label{app:heating_curve}

{%
\captionof{table}{PWL temperature profile (heating curve).}
\label{tab:pwl}
\centering
\footnotesize
\begin{tabular*}{\columnwidth}{@{\extracolsep{\fill}} l c c c c @{}}
\toprule
Load fraction & 0.0 & 0.3 & 0.6 & 1.0 \\
\midrule
$T_{\mathrm{sup}}$ [$^\circ$C] & 70 & 80 & 90 & 100 \\
$T_{\mathrm{ret}}$ [$^\circ$C] & 55 & 50 & 45 & 40 \\
$\Delta T$ [$^\circ$C]         & 15 & 30 & 45 & 60 \\
\bottomrule
\end{tabular*}
}% end centering group

\bigskip

\subsection{L2 zone mapping}
\label{app:l2_zone_mapping}

{%
\captionof{table}{L2 zone mapping (15~junctions $\to$ 7~zones).}
\label{tab:l2_zones}
\centering
\footnotesize
\begin{tabular*}{\columnwidth}{@{\extracolsep{\fill}} l l l l l @{}}
\toprule
Zone & Nodes & Cons. & Assets & Pipe \\
\midrule
A & j$_1$              & V$_1$              & CHP,TES,Gas,Bio,HP,EB & -- \\
B & j$_2$--j$_3$       & V$_2$--V$_3$       & --                    & 650\,m \\
C & j$_4$--j$_5$       & V$_4$--V$_9$       & --                    & 500\,m \\
D & j$_6$--j$_8$       & V$_{10}$--V$_{14}$ & --                    & 450\,m \\
E & j$_9$--j$_{11}$    & V$_{15}$--V$_{18}$ & --                    & 880\,m \\
F & j$_{12}$--j$_{13}$ & V$_{19}$--V$_{24}$ & --                    & 470\,m \\
G & j$_{14}$--j$_{15}$ & V$_{25}$--V$_{27}$ & --                    & 305\,m \\
\bottomrule
\end{tabular*}
}% end centering group

\bigskip

\subsection{Post-processing validation}
\label{app:state_validation}

\noindent\textit{Pipe velocity:}
$v_{p,t} \leq \SI{2.5}{\metre\per\second}$; no violations.
\textit{Pumping (L1/L2/L3):}
${<}5\,$kW peak, ${<}1\,\%$ of thermal losses.
\textit{$\Lnl$ convergence:}
target ${\leq}\,0.5\,\%$ gap; achieved 1.99\,\% at 4\,h time
limit for 744\,h window.

\subsection{Gap stability across scenarios}
\label{app:gap_stability}

{%
\captionof{table}{Gap stability across 11~scenarios (\%).}
\label{tab:gap_stability}
\centering
\footnotesize
\begin{threeparttable}
\begin{tabular*}{\columnwidth}{@{\extracolsep{\fill}} l c c c @{}}
\toprule
Scenario & L1--L3 & L2--L3 & L3--$\Lplus$ \\
\midrule
Baseline     & +13.0 & +2.6 & +0.11\tnote{b} \\
Gas high     & +12.4 & +2.4 & +0.18 \\
Gas low      & +13.5 & +2.8 & +0.05 \\
Elec.\ high  & +12.5 & +2.2 & +0.00 \\
Elec.\ low   & +12.5 & +1.9 & +0.10 \\
CO$_2$ high  & +12.1 & +2.2 & +0.08 \\
CO$_2$ low   & +14.3 & +3.0 & +0.16 \\
Cold year    & --\tnote{a} & +0.6 & +0.11 \\
Warm year    & --\tnote{a} & $-$2.2 & +0.00 \\
HP COP low   & +13.0 & +2.5 & +0.01 \\
Biomass exp. & +13.0 & +2.5 & +0.20 \\
\bottomrule
\end{tabular*}
\begin{tablenotes}\scriptsize
  \item[a] Uniform scaling; L1--L3 not comparable.
  \item[b] Value from primary-case reference run
    (\Cref{tab:cost_extended}); sensitivity scenarios use
    independent solver instances.
\end{tablenotes}
\end{threeparttable}
}% end centering group

%% ============================================================
\section*{CRediT authorship contribution statement}
\textbf{Lukas Ruess:} Conceptualisation, Methodology, Software,
Validation, Formal analysis, Investigation, Data curation,
Writing -- original draft, Visualisation, Funding acquisition.
\textbf{Alexander Sauer:} Supervision, Writing -- review \& editing.

%% ============================================================
\bibliographystyle{elsarticle-num-names}
\bibliography{paper1_dh_fidelity/Paper20_Literatur}
\end{document}